\documentclass[11pt]{article}

\usepackage[T1]{fontenc}
\usepackage[utf8]{inputenc}
\usepackage{lmodern}
\usepackage{microtype}
\usepackage[margin=1in]{geometry}
\usepackage{amsmath,amssymb,amsthm,mathtools,stmaryrd}
\usepackage{booktabs,longtable,tabularx,array}
\usepackage{enumitem}
\usepackage{xcolor}
\usepackage{listings}
\usepackage{fancyhdr}
\usepackage{hyperref}
\usepackage[nameinlink,capitalise,noabbrev]{cleveref}

\definecolor{deepblue}{RGB}{25,63,116}
\definecolor{softgray}{RGB}{245,247,250}
\definecolor{darkgreen}{RGB}{23,96,61}
\definecolor{darkred}{RGB}{145,36,36}

\hypersetup{
  colorlinks=true,
  linkcolor=deepblue,
  citecolor=deepblue,
  urlcolor=deepblue,
  pdftitle={Finitary interfaces and proof-producing compilation for mathematics in iPRA},
  pdfauthor={OpenAI research report for Vladimir Reshetnikov}
}

\pagestyle{fancy}
\fancyhf{}
\fancyhead[L]{Finitary interfaces for iPRA}
\fancyhead[R]{\thepage}
\setlength{\headheight}{14pt}

\setlist[itemize]{leftmargin=*,topsep=3pt,itemsep=2pt}
\setlist[enumerate]{leftmargin=*,topsep=3pt,itemsep=2pt}
\setlength{\parindent}{0pt}
\setlength{\parskip}{5pt plus 1pt minus 1pt}
\emergencystretch=2em

\newtheorem{theorem}{Theorem}[section]
\newtheorem{lemma}[theorem]{Lemma}
\newtheorem{proposition}[theorem]{Proposition}
\newtheorem{corollary}[theorem]{Corollary}
\theoremstyle{definition}
\newtheorem{definition}[theorem]{Definition}
\newtheorem{construction}[theorem]{Construction}
\newtheorem{example}[theorem]{Example}
\theoremstyle{remark}
\newtheorem{remark}[theorem]{Remark}
\newtheorem{warning}[theorem]{Warning}

\newcommand{\N}{\mathbb N}
\newcommand{\Z}{\mathbb Z}
\newcommand{\Q}{\mathbb Q}
\newcommand{\F}{\mathbb F}
\newcommand{\Bool}{\mathbf B}
\newcommand{\iPRA}{\mathrm{iPRA}}
\newcommand{\PRAAx}{\mathsf{PRAAx}}
\newcommand{\Prf}{\mathsf{Prf}}
\newcommand{\PTm}{\mathsf{PTm}}
\newcommand{\PR}{\mathsf{PR}}
\newcommand{\Obs}{\mathsf{Obs}}
\newcommand{\Val}{\mathsf{Val}}
\newcommand{\Good}{\mathsf{Good}}
\newcommand{\Need}{\mathsf{Need}}
\newcommand{\Run}{\mathsf{Run}}
\newcommand{\Inv}{\mathsf{Inv}}
\newcommand{\Accept}{\mathsf{Accept}}
\newcommand{\Cand}{\mathsf{Cand}}
\newcommand{\Edge}{\mathsf{Edge}}
\newcommand{\Viable}{\mathsf{Viable}}
\newcommand{\diam}{\operatorname{diam}}
\newcommand{\htop}{\operatorname{ht}}
\newcommand{\RepoCommit}{\texttt{9df5cbdf39e6}}
\newcommand{\code}[1]{\texttt{\detokenize{#1}}}

\lstdefinelanguage{Lean}{
  morekeywords={def,theorem,lemma,structure,inductive,namespace,end,where,by,exact,have,let,match,with,import,open,variable,forall,if,then,else},
  sensitive=true,
  morecomment=[l]{--},
  morecomment=[s]{/-}{-/},
  morestring=[b]"
}
\lstset{
  basicstyle=\ttfamily\small,
  backgroundcolor=\color{softgray},
  frame=single,
  rulecolor=\color{black!15},
  breaklines=true,
  columns=fullflexible,
  keepspaces=true,
  showstringspaces=false
}

\title{\textbf{Finitary Interfaces and Proof-Producing Compilation\\for Mathematics in Intuitionistic Primitive Recursive Arithmetic}\\[0.6em]
\large Finite observations, relational certificates, guarded operations, and an object-language trace compiler}
\author{Research report prepared for Vladimir Reshetnikov}
\date{15 November 2026\\Repository snapshot: \RepoCommit}

\begin{document}
\maketitle

\begin{center}
\fbox{\begin{minipage}{0.91\textwidth}
\textbf{Exact status.} The mathematical arguments in this report are human-readable proofs and proof schemas. The included Python reference layer has been executed with exact rational arithmetic. The Lean files are an uncompiled, source-oriented outline against the pinned ProveIt API and contain visible \code{sorry} obligations. No new kernel-checked theorem is claimed. The already checked bridge \code{pra_iff_native} is inherited from \code{ConciseIPRA} and applies after an ordinary \(\Prf\,\PRAAx\) derivation has been completed.
\end{minipage}}
\end{center}

\begin{abstract}
The ProveIt iPRA tree contains many finitary reconstructions of classical mathematics and several cross-domain frameworks based on certificates, finite observations, explicit moduli, and relational witnesses. Their common difficulty is not the mere coding of finite mathematics. It is the simultaneous control of four boundaries: uniform primitive recursiveness, constructive choice, quantitative adequacy, and production of an actual object-language derivation of type \(\Prf\,\PRAAx\).

This report develops a unified \emph{finitary interface calculus}. A mathematical object is represented only through primitive-recursive finite observations. A construction is a contracted transducer with an explicit backward precision demand, a bounded verifier, optional finite candidate clouds, and local object-language proofs. Fixed bounded formulas are Booleanized into quantifier-free equations. Bounded computations are internalized by a finite-trace reflection theorem using the existing \code{PRA.qfInd}. Contract DAGs are compiled by a reverse demand pass and one forward trace. Nonunique approximate witnesses are handled by backward viability and stable quotient erasure, not by hidden global selectors. Infinite objects arise only after a diameter or tail modulus proves that finite observations cohere.

The framework resolves, at the level of complete mathematical arguments, the recurring issues of all-precision induction, partial operations, branch discontinuity, finite compactness, quotient extensionality, and the finite-list/height-bound interface. It also isolates genuine residual obligations: general height producers, rational covering representatives, ineffective analytic constants, and unrestricted exact branch selection. Detailed Lean definitions and a dependency-ordered path to a complete formalization are supplied, together with an exact Python reference kernel and reproducible certificates.
\end{abstract}

\tableofcontents
\newpage

\section{Executive synthesis}

\subsection{The problem is an interface problem}

A proof in intuitionistic primitive recursive arithmetic is not obtained merely by observing that a familiar argument is finite ``in principle.'' Four distinct claims must be separated.

\begin{enumerate}
\item \textbf{Representability.} Every object manipulated by the argument must have a fixed natural-number coding, and every operation used by the argument must belong to a fixed primitive-recursive family.
\item \textbf{Quantitative adequacy.} Whenever the ordinary theorem invokes convergence, compactness, continuity, a finite cover, an isolated branch, or a height bound, the formal statement must expose the numerical modulus or certificate actually consumed by the construction.
\item \textbf{Constructive coherence.} Finite witnesses obtained at different precisions need not agree. One must either prove that all valid choices have the same observable effect or build a compatible finite path by bounded search. A hidden appeal to an infinite branch, an exact sign, or a canonical label is not acceptable.
\item \textbf{Object-language reflection.} A correct Lean function, Python program, or external finite proof does not by itself yield an iPRA theorem. The end product must be an actual derivation of type
\[
   \Gamma \;\vdash[\PRAAx]\; \varphi,
\]
that is, a value of the repository's inductive relation \(\Prf\,\PRAAx\).
\end{enumerate}

The repository already addresses the first three points in many individual packets and several general frameworks. The most persistent cross-packet gap is the fourth. The central contribution of the present package is therefore a proof-producing interface whose semantic and syntactic sides meet at a small, explicit theorem: finite-trace reflection by quantifier-free induction.

\subsection{The five-layer architecture}

Every formalization target is divided into five layers. The division is not cosmetic: it is the main safeguard against accidentally treating a host-language theorem as an object-theory proof.

\paragraph{Layer I: fixed presentation.}
A presentation is a meta-level package of fixed primitive-recursive symbols. It may describe finite data exactly or infinite data by finite observations. It is deliberately not a natural-number code for an arbitrary program. This avoids the impossible demand for a total primitive-recursive universal evaluator.

\paragraph{Layer II: quantitative contract.}
A construction advertises the precision, search, guard, and certificate data it needs. Its contract consists of a primitive-recursive demand map, a bounded verifier, and an observable postcondition. A producer may be included, but a verifier is meaningful even when the difficult producer remains an open mathematical obligation.

\paragraph{Layer III: finite execution.}
The construction is executed as a bounded trace, a finite contract DAG, a bounded search, or a finite relational dynamic program. Every loop has an explicit primitive-recursive bound; every branch is Boolean or remains relational; every finite quantifier has a displayed bound.

\paragraph{Layer IV: object proof.}
The local base, step, and final contracts are themselves \(\Prf\,\PRAAx\) derivations. A generic compiler combines them using substitution, cut, the defining equations of primitive recursion, and \code{PRA.qfInd}. This layer contains no appeal to the intended standard model.

\paragraph{Layer V: transport.}
Only after the ordinary iPRA derivation is complete is it transported to the small native scalar core by the already checked theorem \code{pra_iff_native}. Keeping transport last makes the trust boundary visible.

\subsection{A single normal form for many proofs}

A large class of packet arguments can be reorganized into the following statement. Fix input \(x\), requested precision \(q\), and a bounded state evolution
\[
  s_0=I(x,q),\qquad s_{i+1}=S(i,s_i,x,q),\qquad i<N(x,q).
\]
Let \(\Phi(i,s,x,q)\) be a quantifier-free invariant after bounded-formula compilation. Suppose iPRA proves the three local obligations
\begin{align*}
  \Good(x)&\to \Phi(0,I(x,q),x,q),\\
  \Good(x)\wedge \Phi(i,s,x,q)&\to \Phi(i+1,S(i,s,x,q),x,q),\\
  \Good(x)\wedge \Phi(N(x,q),s,x,q)&\to \Accept(s,x,q).
\end{align*}
Then the primitive-recursive run symbol \(\Run\) satisfies
\[
  \iPRA\vdash \Good(x)\to
  \Accept(\Run(N(x,q),x,q),x,q).
\]
The proof is a direct application of quantifier-free induction to
\(
  \Phi(i,\Run(i,x,q),x,q).
\)
The recursion equations rewrite the zero and successor cases. This is elementary mathematics, but it is exactly the elementary mathematics that turns a bounded algorithm into an object-language proof.

The same normal form handles matrix elimination, finite search, interval bisection, finite-net minimization, construction of a simplicial chain map, bounded descent, certificate checking, and the forward pass of a contract DAG. A complicated theorem may still require difficult estimates, but once those estimates are presented as finite contracts, the proof-engineering problem becomes uniform.

\subsection{What is genuinely resolved}

The framework gives complete non-formalized proofs of the following reusable facts.

\begin{itemize}
\item Every \emph{fixed} bounded formula over primitive-recursive atomic predicates has a primitive-recursive Boolean characteristic function. Its truth is equivalent in iPRA to the quantifier-free equation that the characteristic function equals one.
\item External parameters, including precision, should be fixed before induction and universally quantified afterward. This avoids putting a quantifier over all precisions inside an induction predicate.
\item A finite contract DAG has a sound reverse demand pass. At fan-out, demands combine by maximum. A single forward bounded trace can check all local contracts.
\item Backward viability on finite candidate clouds is equivalent to the existence of a complete compatible finite path. Choosing the least viable successor is primitive recursive and invokes no infinite dependent choice.
\item If the terminal observable is constant on all viable paths, or has an explicit diameter bound, internal tie-breaking can be erased. This is the right replacement for a discontinuous canonical selector.
\item A partial numerical operation becomes an iPRA construction after totalization and the addition of a positive rational guard. Mere negated equality is not enough for reciprocal, pivoting, chart selection, or branch isolation.
\item Effective compactness is represented by finite nets, a coverage certificate, and a modulus. Approximate minimization is then bounded search.
\item For a decidable coded-point predicate, an explicit height bound yields a complete finite list, and a certified complete finite list yields a height bound by finite maximum. Thus the deep and elementary parts of Diophantine finiteness separate cleanly.
\end{itemize}

\subsection{What remains honestly open}

The framework does not manufacture data that the underlying mathematics does not supply. In particular it does not prove a general effective Faltings or Siegel height bound; it does not turn Selmer classes into rational points on covering spaces; it does not make ineffective constants effective; it does not uniformly choose a sign of an arbitrary real name, a Puiseux branch from arbitrary Cauchy coefficients, or a rotation axis at the identity; and it does not turn metastability into a Cauchy modulus.

These are not implementation omissions of the same kind as a missing substitution lemma. They are producer obligations or computability barriers. The framework's value is that it makes the distinction machine-visible. A packet can now be classified as having
\begin{enumerate}
\item a complete finite verifier,
\item an explicit but unproved producer theorem,
\item a relational output that avoids the producer, or
\item an impossible unrestricted specification that must be reformulated.
\end{enumerate}

\subsection{Deliverables in this archive}

The archive contains:
\begin{itemize}
\item this report and its complete \LaTeX{} source;
\item Lean source blueprints using the exact pinned \code{ConciseIPRA} API;
\item a dependency-ordered Lean completion plan;
\item an exact Python reference kernel implementing bounded folds, trace checking, contract-DAG demand propagation, finite viability, three-way guards, guarded reciprocal, finite-net minimization, height/list conversion, and cubic bisection;
\item eleven passing regression tests and a serialized cubic-inverse certificate;
\item a theorem ledger, a difficulty matrix, a machine-readable framework schema, source provenance, hashes, and PDF validation records.
\end{itemize}

The Python layer is deliberately modest. Its purpose is to expose finite algorithms and catch quantifier-order or certificate-format mistakes. It is not counted as a proof of an iPRA theorem. The Lean files are similarly explicit about their status: they give definitions, exact theorem signatures, and small remaining milestones, but retain visible \code{sorry} placeholders and were not compiled in the build environment.

\section{Repository audit and the precise missing bridge}

\subsection{Pinned source tree}

This report was prepared against the public ProveIt repository at commit
\RepoCommit.  The relevant subtree distinguishes two kinds of material.
The \code{ConciseIPRA} directory is a Lean development whose kernel-checked
results include the ordinary first-order presentation of iPRA, a compact native
scalar calculus, and translations between them.  The neighboring
\code{docs/proofs} directory is a research archive.  At the pinned snapshot its
index records seventy-two delivered packets in algebra, analysis, analytic and
Diophantine number theory, combinatorics, foundations, geometry, probability,
and topology.  The index explicitly warns that these packets were filed as
delivered: their Lean files were not systematically built, and their claims
must not be confused with the checked core \cite{ProveItRoot,PacketsIndex}.

This distinction determines the standard of evidence used here.  A theorem in
this report is called \emph{embedded} only when there is a specified route to a
value of the repository's inductive derivability relation
\[
        \Prf\,\PRAAx\,\Gamma\,\varphi.
\]
A host-language proof about natural numbers, a successful exact-arithmetic
program, or a semantic theorem about a Lean structure is useful evidence, but
it is not itself such a value.

\subsection{What the existing frameworks already accomplish}

The foundation packets are unusually convergent.  Although they use different
terminology, they repeatedly discover the same useful replacements for
classical infinitary objects.

\begin{longtable}{p{0.21\textwidth}p{0.34\textwidth}p{0.37\textwidth}}
\toprule
Classical idiom & Finitary replacement & Recurring repository contribution \\
\midrule
\endfirsthead
\toprule
Classical idiom & Finitary replacement & Recurring repository contribution \\
\midrule
\endhead
A completed infinite object & finite observations indexed by precision &
certified finitary presentations, observation towers, finite names \\
A canonical approximate witness & a finite cloud of verified answers &
relational certificates, root clusters, stable observables \\
Continuity ``by compactness'' & an explicit modulus plus a finite net &
finite observation and uniform-certificate frameworks \\
A long algorithmic argument & a bounded trace with a quantifier-free invariant &
finite-demand, local--global, and compiler packets \\
Nested numerical constructions & backward propagation of required precision &
finite-demand and contract-DAG architectures \\
A partial operation & a total program plus a positive rational guard &
guarded inverses, pivots, charts, and branch rules \\
Finiteness of a Diophantine set & a height producer plus bounded enumeration &
height/list interfaces and finite witness packages \\
Equality of abstract objects & observational equivalence at every finite level &
setoid and relational back-and-forth methods \\
\bottomrule
\end{longtable}

The most developed unified packet proves human-level theorems for bounded
Booleanization, finite trace reflection, relational viability, shrinking
candidate clouds, guarded arithmetic, finite nets, and height/list conversion.
Its theorem ledger nevertheless marks the generic \code{Prf PRAAx} compiler as
an open engineering task.  The subsequent contract-DAG packet strengthens the
architecture with reverse demand propagation and a global trace, but again
leaves the generic object-language theorem unimplemented.  This is the point
at which the present work begins: it treats those semantic ideas as prior art
and tightens the exact interface with the checked \code{ConciseIPRA} API.

\subsection{The checked backend that can be reused}

The file \code{PR.lean} supplies precisely the syntax needed for the target.
It defines primitive-recursive symbols
\[
 \PR(k),
\]
terms \(\PTm(n)\), the axiom family \(\PRAAx\), and the derived rule
\code{PRA.qfInd}.  Primitive recursion is represented by
\code{PR.prec}; its zero and successor equations are available as theorems
\code{PRA.precZero} and \code{PRA.precSucc}.  The generic deduction file
provides substitution, weakening, simultaneous cut, implication and universal
rules, and Leibniz transport.  In particular, the following operations are
already checked:
\[
\begin{array}{ll}
\code{Prf.subst}, & \text{substitute terms for all free variables},\\
\code{Prf.weaken},& \text{enlarge a context},\\
\code{Prf.trans}, & \text{replace hypotheses by their derivations},\\
\code{Prf.eqSubst1},& \text{transport a formula along an equality},\\
\code{PRA.qfInd},& \text{induct on a quantifier-free formula}.
\end{array}
\]
Finally, \code{NativeBiInterpretation.lean} proves
\code{pra_iff_native}: an ordinary iPRA derivation is equivalent to a
derivation of its reading in the native scalar calculus.  Consequently no new
bi-interpretation theorem is needed.  The missing work is entirely on the
ordinary \(\Prf\,\PRAAx\) side.

\subsection{Classification of the reported difficulties}

The difficulty reports in the packets can be separated into four classes.
This classification is essential because only the first class can be solved
by ordinary Lean engineering.

\begin{definition}[Four-way status classification]
A target or subtarget has one of the following statuses.
\begin{enumerate}
\item \emph{Compiler obligation}: all mathematical data and local proofs are
available, but the generic construction of an object-language derivation has
not been implemented.
\item \emph{Quantitative producer obligation}: a finite verifier is known,
but an explicit bound, modulus, representative, or finite cover remains to be
proved.
\item \emph{Representation defect}: the requested output is discontinuous or
noncanonical; a relational or guarded statement is required.
\item \emph{Computability barrier}: the unrestricted producer would imply a
nonconstructive principle or solve an undecidable problem.
\end{enumerate}
\end{definition}

Examples of the first class are bounded-fold reflection and finite-trace
induction.  Examples of the second are general Siegel/Faltings height bounds,
representatives on Mordell--Weil covering curves, and effective constants in
certain prime-distribution results.  Exact root labels at collisions and an
axis for the identity rotation are representation defects.  A total primitive-
recursive universal evaluator and a uniform conversion from mere finiteness to
a search bound are barriers.

The framework below solves the compiler obligations, gives a common interface
for producer obligations, repairs representation defects, and records barriers
rather than concealing them behind axioms.

\subsection{Audit conclusion}

The repository does not need another informal assertion that ``all the
relevant computations are finite.''  It needs a small compilation theorem
whose hypotheses are already the kind of local formulas the packets can
produce and whose conclusion is an actual iPRA derivation.  Once that theorem
is in place, the many domain-specific packets can share a single conversion
workflow:
\[
\begin{aligned}
\text{mathematical proof}
&\longrightarrow \text{finite contracted program}\\
&\longrightarrow \text{local }\Prf\,\PRAAx\text{ obligations}\\
&\longrightarrow \text{global }\Prf\,\PRAAx\text{ derivation}.
\end{aligned}
\]
The remaining sections develop this pipeline and prove its mathematical
correctness.

\section{Logical target and standards of formalization}

\subsection{iPRA as used in the repository}

For this report, iPRA is first-order intuitionistic logic with equality over a
function symbol for every primitive-recursive construction, the defining
equations for those symbols, the usual successor separation and injectivity
axioms, and induction for quantifier-free formulas.  The checked repository
implementation makes this precise:
\begin{itemize}
\item \(\PR(k)\) is the inductive type of primitive-recursive symbols of arity
\(k\), generated by zero, successor, projections, composition, and primitive
recursion;
\item \(\PTm(n)\) is the type of terms with at most \(n\) free de Bruijn
variables;
\item \(\PRAAx\) is the family of open iPRA axioms;
\item \(\Prf\,\PRAAx\,\Gamma\,\varphi\) is intuitionistic derivability from
context \(\Gamma\).
\end{itemize}

The restriction to quantifier-free induction is not a serious obstacle to
finite computation.  It is a design constraint: state, loop counters,
certificates, and all bounded quantifiers must be compiled into terms and
Boolean equations before induction is invoked.

\subsection{Three levels that must not be conflated}

A robust formalization keeps three levels visible.

\paragraph{Meta-level specification.}
Lean structures may organize arities, formulas, and proof terms.  For example,
a structure may contain fixed symbols \(\Good,\Run,\Inv,\Accept\) and fields
whose types are local \(\Prf\,\PRAAx\) derivations.  This is harmless: Lean is
only assembling an explicit object-language proof.

\paragraph{Host-level semantic validation.}
An executable program can test examples, discover missing bounds, and validate
finite certificate formats.  Such tests establish that the proposed finite
algorithm behaves as intended in the standard model.  They do not establish
that iPRA proves its correctness.

\paragraph{Object-language derivation.}
The final theorem is a constructor returning \(\Prf\,\PRAAx\).  Every use of
induction, equality transport, bounded search, or a local theorem is visible in
that constructor.  Only this level is transported by \code{pra_iff_native}.

\begin{warning}[Semantic reflection is not enough]
A theorem of the form
\[
  \forall u\in\N,\quad \operatorname{run}(u)\text{ satisfies }\Accept
\]
proved in Lean's ambient logic is not interchangeable with
\[
  []\vdash[\PRAAx]\forall u\,\Accept(\Run(\operatorname{bound}(u),u),u).
\]
The former may use unrestricted induction, higher-order choice, or library
facts unavailable to iPRA.  The framework requires local object-language
premises and constructs the latter directly.
\end{warning}

\subsection{Fixed families rather than a universal evaluator}

There is no total primitive-recursive universal evaluator for all primitive-
recursive functions.  Therefore a valid framework cannot quantify over an
arbitrary code \(e\) and assume a total PR term \(U(e,x)\) evaluating it.
There are three safe alternatives.

\begin{enumerate}
\item \textbf{Fixed symbols.}  For each formalized theorem, choose finitely many
symbols in \(\PR(k)\).  This is the simplest and strongest route.
\item \textbf{Fixed program family.}  Define an inductive family of programs
whose interpreter is primitive recursive because a structural size or a
syntactic recursion depth is fixed by the family.
\item \textbf{Clocked computation.}  Evaluate a code for at most \(t\) steps.
The bounded evaluator is primitive recursive; a trace certificate supplies
\(t\) and proves acceptance.
\end{enumerate}

The framework is parametric at the Lean level but concrete in the object
language.  A theorem may accept a structure containing arbitrary \emph{fixed}
\(\PR\) symbols.  It never asserts the existence of one internal evaluator for
all such structures.

\subsection{The target shape of a reusable theorem}

The basic compiler theorem has the following logical form.  Let
\(G(u)\), \(I(i,s,u)\), and \(A(s,u)\) be quantifier-free predicates represented
by Boolean equations, and let \(r(i,u)\) be the primitive-recursive run.
Assume closed derivations of
\begin{align*}
 &G(u)\to I(0,r(0,u),u), \\
 &G(u)\to I(i,s,u)\to I(Si,\operatorname{step}(i,s,u),u),\\
 &G(u)\to I(b(u),s,u)\to A(s,u).
\end{align*}
Then the compiler must return a closed derivation of
\[
   \forall u\,[G(u)\to A(r(b(u),u),u)].
\]
The theorem is deliberately intensional.  Its inputs are not claims that the
three implications are true in \(\N\); they are values of the corresponding
\(\Prf\,\PRAAx\) types.

\subsection{What counts as a complete Lean outline}

A useful outline must expose enough information that completion is routine
proof engineering rather than renewed mathematical research.  In this package
that means:
\begin{itemize}
\item every term builder has a fixed arity and a documented de Bruijn layout;
\item every formula used by induction is explicitly defined;
\item the quantifier-free side condition is a named lemma;
\item local hypotheses and the global result have exact
\(\Prf\,\PRAAx\) types;
\item the zero and successor run equations are tied directly to
\code{PRA.precZero} and \code{PRA.precSucc};
\item the remaining proof is split into base, step, induction, and final
milestones, each using named checked APIs;
\item every unproved part remains marked by \code{sorry}; no compilation claim
is made.
\end{itemize}

This is stronger than a pseudocode sketch but weaker than a completed kernel
proof.  It gives a clear path to completion and makes the residual work
measurable.

\section{Primitive-recursive coding library}

\subsection{Why a common coding layer matters}

Most ordinary mathematical objects used in the packets---finite lists,
matrices, rational numbers, polynomials, simplicial complexes, proof traces,
and finite covers---are easily coded by natural numbers.  The danger is not
existence of a coding.  It is silent dependence on unbounded decoding,
ill-formed inputs, or quotient representatives.  A common coding layer should
therefore satisfy three properties:
\begin{enumerate}
\item all constructors and bounded destructors are primitive recursive;
\item every algorithm is total on all natural-number inputs;
\item correctness theorems are guarded by a primitive-recursive well-formedness
predicate.
\end{enumerate}

\subsection{Sequences and bounded tables}

Fix a primitive-recursive pairing function
\(\langle x,y\rangle\) with primitive-recursive projections
\(\pi_0,\pi_1\).  A finite sequence may be coded by nested pairing, by a prime-
power code, or by a bounded binary block.  The exact choice is irrelevant to
proof strength; the following abstract interface is sufficient:
\[
\operatorname{len}(a),\qquad
\operatorname{get}(a,i),\qquad
\operatorname{set}(a,i,x),\qquad
\operatorname{push}(a,x).
\]
All operations are total.  Out-of-range reads return zero, and out-of-range
writes leave the code unchanged.  The predicate \(\operatorname{Seq}(a)\)
certifies canonical shape when canonicality matters.

A rectangular table is represented by a triple \((r,c,a)\), with the entry at
\((i,j)\) stored at offset \(ic+j\).  Matrix addition, multiplication, minors,
and determinants of a fixed finite dimension are primitive recursive.  When
dimension is variable, the algorithms remain primitive recursive because
loops are bounded by the coded dimensions.

\subsection{Integers and rationals}

Represent an integer by a pair \((p,n)\), denoting \(p-n\).  This avoids a
case distinction in subtraction.  Equality and order reduce to natural-number
arithmetic:
\[
  (p,n)=(p',n')\quad\Longleftrightarrow\quad p+n'=p'+n.
\]
A rational code is a triple \((p,n,d)\) with \(d>0\), denoting
\((p-n)/d\).  Algorithms are total even at \(d=0\), using a default value;
well-formedness supplies the positive-denominator guard.  Rational equality
uses cross multiplication and therefore does not require normalization.

For uniqueness or finite enumeration, one may add the canonical predicate
\[
  d>0,\qquad \gcd(|p-n|,d)=1,
\]
with a fixed sign convention.  Canonicalization is primitive recursive by the
Euclidean algorithm.  The framework does not force canonicalization in every
inner loop; it uses observational equality where that is cheaper.

\subsection{Finite algebraic data}

A polynomial over a coded ring is a coefficient sequence.  Evaluation,
derivative, truncated multiplication, Euclidean pseudo-division, resultants,
and determinant-based certificates are bounded recursions on degree.
A finitely presented algebra is encoded by:
\begin{itemize}
\item numbers of generators and relations;
\item coefficient tables for the relations;
\item a degree or monomial cap when a finite normal-form computation is used;
\item a certificate that all reductions stay below that cap.
\end{itemize}

The cap is important.  ``Run Buchberger's algorithm until it terminates'' is
not directly a primitive-recursive program unless a bound or a decreasing
measure is supplied.  A finite Macaulay matrix at an explicitly given degree,
by contrast, is a bounded object and its row-reduction certificate is directly
verifiable.

\subsection{Finite combinatorial and topological data}

A finite graph is a vertex bound together with a Boolean edge table.  A finite
simplicial complex is a vertex bound and a Boolean predicate on bounded finite
subsets, with a downward-closure certificate.  Chains over \(\F_2\) are bit
vectors; boundary maps are finite matrices.  Homology statements at a fixed
finite level can be expressed using matrix equations:
\[
  \partial z=0,
  \qquad z=\partial w,
  \qquad f_\#\partial=\partial f_\#.
\]
No quotient type must be internalized.  ``Same homology class'' is a bounded
existential statement about a chain \(w\), which can be Booleanized.

Finite categories are coded by bounds on objects and arrows, source/target and
composition tables, and Boolean laws for identity and associativity.  A
functor is a pair of tables satisfying bounded equations; a natural
transformation is another finite table.  Thus finite limits, finite colimits,
and finite Yoneda calculations reduce to bounded search and equality.

\subsection{Traces and certificates}

A proof trace has a fixed state encoding.  Its code need not expose a list of
all states: a primitive-recursive run term is often smaller.  For external
certificate validation, however, a list form is convenient.  We use both:
\begin{itemize}
\item the object language defines \(\Run(i,u)\) by primitive recursion;
\item a JSON or host-level certificate may list states
\(s_0,\ldots,s_b\) and is checked against the same local step relation.
\end{itemize}

The trace invariant must be quantifier free.  Bounded subcertificates are
therefore folded into a Boolean state component rather than quantified over in
the induction predicate.

\subsection{Observational setoids}

Many mathematical objects have several codes.  Instead of forcing a global
canonical form, define at each precision \(p\) a primitive-recursive relation
\[
      a\equiv_p b
      \quad\Longleftrightarrow\quad
      \Obs(a,p)=\Obs(b,p)
\]
or, more generally, a bounded back-and-forth relation between finite answer
clouds.  The required laws are finite:
\begin{enumerate}
\item reflexivity, symmetry, and transitivity at each precision;
\item monotonicity: equivalence at a stronger precision implies equivalence at
weaker precision;
\item congruence of every contracted operation.
\end{enumerate}

This setoid discipline resolves the frequent confusion between equality of
codes and equality of represented objects.  It also permits efficient internal
representations while preserving extensional theorem statements.

\subsection{Coding theorem}

\begin{theorem}[Bounded data are iPRA-ready]
Fix a finite signature of sequence, table, integer, rational, graph, matrix, or
finite-complex operations built from the interfaces above.  Every constructor,
bounded destructor, well-formedness test, and bounded verifier is primitive
recursive.  Every fixed equality or validity condition is representable by a
quantifier-free Boolean equation in iPRA.
\end{theorem}

\begin{proof}
Pairing and its projections are primitive recursive.  The listed operations
are obtained by composition, bounded iteration, and primitive recursion on an
explicit natural bound.  Boolean combinations preserve primitive
recursiveness.  Each Boolean characteristic function is named by a symbol in
\(\PR\); the equation to the Boolean value one is quantifier free.  No
unbounded minimization is used.  Total default behavior handles malformed
codes, while correctness theorems assume the corresponding well-formedness
Boolean.
\end{proof}

\section{Certified finitary presentations}

\subsection{Deterministic presentations}

A finitary presentation describes only the observations of an object that a
finite proof can consume.

\begin{definition}[Deterministic finitary presentation]
A deterministic finitary presentation \(X\) consists of fixed primitive-
recursive data:
\begin{enumerate}
\item a well-formedness predicate \(\Good_X(x)\);
\item an observation function \(\Obs_X(x,p)\), where \(p\) is a precision or
finite query code;
\item a restriction function \(\operatorname{res}_X(p,q,a)\) for weakening an
answer at precision \(p\) to a weaker precision \(q\);
\item a primitive-recursive relation \(q\preceq_X p\);
\item object-language proofs of identity and composition of restriction and of
observation coherence.
\end{enumerate}
\end{definition}

The precision order need not be ordinary numerical order.  It may be a tuple
of a radius, derivative order, denominator cap, finite open cover, degree cap,
or a finite list of requested observables.  All that matters is that its
comparison and joins are primitive recursive.

\begin{example}[Regular real name]
A regular rational name has observations \(q_p\in\Q\) satisfying
\(|q_p-q_{p+1}|\le 2^{-p}\).  Restriction simply discards excess precision.
The presentation need not contain a completed real number; every theorem is a
finite statement about rational approximations and an explicit error budget.
\end{example}

\subsection{Relational presentations}

Deterministic observations are too strong when valid finite witnesses are not
unique.  Root orderings, triangulations, branch choices, local charts, and
finite covers may change discontinuously.

\begin{definition}[Finitary relational presentation]
A relational presentation \(X\) consists of:
\begin{enumerate}
\item a bound \(B_X(x,p)\) on candidate codes;
\item a decidable validity predicate \(\Cand_X(x,p,a)\) for
\(a<B_X(x,p)\);
\item a compatibility relation
\(\Edge_X(x,p,q,a,b)\) between candidates at related precisions;
\item a finite observable \(\Val_X(x,p,a)\);
\item optional diameter and motion bounds controlling variation among valid
or compatible candidates.
\end{enumerate}
\end{definition}

A candidate cloud is therefore a finite relation, not a set-valued object in a
higher-type universe.  Existence means a bounded disjunction.  Compatibility
means a Boolean table.  Backward viability, developed in \cref{sec:relational},
turns these finite tables into coherent finite paths.

\subsection{Contracted transducers}

\begin{definition}[Contracted transducer]
A contracted transducer \(F:X\to Y\) consists of:
\begin{enumerate}
\item a primitive-recursive backward demand
\(\Need_F(x,p_Y)\), specifying the input precision needed for output precision
\(p_Y\);
\item a bounded state transition or finite relation producing output
candidates from valid input candidates;
\item a local verifier and a quantifier-free invariant;
\item an object-language proof that valid input observations at the demanded
precision yield a valid output observation at \(p_Y\);
\item an extensionality proof showing that observationally equivalent inputs
produce equivalent outputs.
\end{enumerate}
\end{definition}

This definition makes precision flow contravariant and data flow covariant.
The reverse demand is part of the theorem, not an informal estimate hidden in
its proof.

\subsection{Composition}

Suppose \(F:X\to Y\) and \(G:Y\to Z\) are deterministic contracted
transducers.  In the simplest, \emph{static-demand} case the demand of \(G\)
depends only on the requested output query (and fixed finite metadata), so
\[
  \Need_{G\circ F}(x,p_Z)
    =\Need_F\bigl(x,\Need_G(p_Z)\bigr).
\]
More generally, \(G\)'s demand may depend on a designated finite summary of its
input.  Then the contract must also state a primitive-recursive preliminary
demand sufficient to compute that summary, followed by a certified upper
envelope for the final demand.  Equivalently, one may join the demands of all
finitely possible guarded branches.  It is not legitimate to write
\(\Need_G(F(x),p_Z)\) if computing that demand already requires an unbounded or
unknown-precision observation of \(F(x)\).  For a finite program DAG,
\cref{sec:dag} packages the static or majorized case into a reverse pass and a
single forward trace.

\begin{theorem}[Sequential contract composition]
The composite of two contracted transducers is a contracted transducer.  Its
local soundness derivation is obtained by substitution and simultaneous cut
from the local derivations of the factors.
\end{theorem}

\begin{proof}
Fix an output query \(p_Z\).  The contract of \(G\) produces a finite demand
on \(Y\); the contract of \(F\) produces the corresponding demand on \(X\).
Run \(F\) and then \(G\).  The first local derivation concludes the hypotheses
required by the second.  In the repository calculus, instantiate both closed
local proofs with \code{Prf.subst}, use \code{Prf.trans} to replace the
intermediate assumptions, and discharge the input contract by implication
introduction.  Since all formula builders are fixed and quantifier free, the
composite verifier remains quantifier free.
\end{proof}

\subsection{Products, sums, and finite families}

Products combine demands by componentwise join.  A finite family of consumers
of the same producer combines demands by maximum or by a primitive-recursive
join operation in the query preorder.  Tagged sums are safe when the tag is
part of the finite input.  When a tag is selected by approximate numerical
data, one must use a three-way guard or a relational output rather than exact
binary branching.

\subsection{Proof-carrying presentations}

A presentation can carry two distinct kinds of proof.

\begin{enumerate}
\item A \emph{certificate verifier} is a fixed Boolean program.  It checks a
finite witness such as a factorization identity, a matrix inverse, a cover, or
a homotopy.
\item A \emph{proof-producing compiler} is a Lean function whose input fields
are local \(\Prf\,\PRAAx\) derivations and whose output is a global derivation.
\end{enumerate}

The first kind makes domain mathematics modular: a deep theorem may be
isolated as a producer of a finite certificate.  The second kind prevents the
framework itself from becoming merely semantic.

\subsection{Adequacy and honesty}

A presentation is adequate for a theorem only if its observations determine
the theorem's conclusion at the claimed precision.  This requirement rules
out several common shortcuts:
\begin{itemize}
\item a metastable window is not a Cauchy name without a tail modulus;
\item a finite cardinality bound is not an enumerating cutoff;
\item a nonzero proposition is not a numerical inverse guard;
\item a multiset of roots is not an ordered list at a collision;
\item a Selmer class is not a rational point on its covering curve.
\end{itemize}

The framework does not insist on stronger data than the theorem needs.  It
instead insists that every consumed datum be named and verified.  This is the
main sense in which the interface is both constructive and reusable.

\section{Bounded-formula Booleanization}
\label{sec:booleanization}

\subsection{The bounded language}

Let atomic predicates be equations or fixed primitive-recursive Boolean tests.
A bounded formula is generated by
\[
\begin{aligned}
\phi ::= {}& P(\vec x)\mid \bot\mid
(\phi\wedge\psi)\mid(\phi\vee\psi)\mid(\phi\to\psi)\\
&\mid \forall y<t(\vec x)\,\phi(\vec x,y)
\mid \exists y<t(\vec x)\,\phi(\vec x,y).
\end{aligned}
\]
The formula is \emph{fixed}: its syntax tree does not arrive as an arbitrary
runtime code.  Parameters and bounds may vary, but the compiler is generated
once for this particular tree.

Use Boolean values \(0,1\), with total operations
\[
  \neg b=1-b,\qquad
  b\wedge c=bc,
  \qquad
  b\vee c=\min(1,b+c),
  \qquad
  b\Rightarrow c=(\neg b)\vee c.
\]
All are primitive recursive.

\subsection{Bounded folds}

Define
\begin{align*}
 \operatorname{All}(0,f)&=1,
 &\operatorname{All}(Sn,f)&=
    \operatorname{All}(n,f)\wedge f(n),\\
 \operatorname{Any}(0,f)&=0,
 &\operatorname{Any}(Sn,f)&=
    \operatorname{Any}(n,f)\vee f(n).
\end{align*}
For each fixed function body \(f\), these are primitive-recursive symbols.
There is no internal higher-type argument: the displayed notation abbreviates
a symbol constructed by primitive recursion from the fixed body.

The least bounded witness is likewise primitive recursive:
\[
 \operatorname{least}(n,f)=
 \begin{cases}
  \min\{i<n:f(i)=1\},&\text{if such }i\text{ exists},\\
  n,&\text{otherwise}.
 \end{cases}
\]
It may be implemented by a scan carrying a ``found'' bit and the current
candidate.

\begin{lemma}[Fold reflection]
For every fixed Boolean term \(f(i,\vec x)\), iPRA proves
\begin{align*}
 \operatorname{All}(n,f)=1
 &\leftrightarrow \forall i<n\,[f(i,\vec x)=1],\\
 \operatorname{Any}(n,f)=1
 &\leftrightarrow \exists i<n\,[f(i,\vec x)=1].
\end{align*}
Moreover, if \(\operatorname{Any}(n,f)=1\), then
\(\operatorname{least}(n,f)<n\) and
\(f(\operatorname{least}(n,f),\vec x)=1\).
\end{lemma}

\begin{proof}
Induct on \(n\).  The zero cases reduce to the defining equations and the
empty bounded quantifiers.  In the successor case, unfold one fold equation.
For \(\operatorname{All}\), Boolean conjunction equals one exactly when both
arguments equal one; apply the induction hypothesis to the prefix and join the
last instance.  The \(\operatorname{Any}\) proof is dual.  For least search,
strengthen the induction predicate to include: either no witness has occurred
and the state equals the current bound, or the stored index is the least
witness seen so far.  This predicate is a fixed Boolean combination of
equalities and inequalities, hence quantifier free after its prefix conditions
are folded.
\end{proof}

\subsection{Structural Booleanization}

Define a characteristic term \(\llbracket\phi\rrbracket\) recursively:
\[
\begin{array}{rcl}
\llbracket P\rrbracket &=& \chi_P,\\
\llbracket\bot\rrbracket&=&0,\\
\llbracket\phi\wedge\psi\rrbracket&=&
 \llbracket\phi\rrbracket\wedge\llbracket\psi\rrbracket,\\
\llbracket\phi\vee\psi\rrbracket&=&
 \llbracket\phi\rrbracket\vee\llbracket\psi\rrbracket,\\
\llbracket\phi\to\psi\rrbracket&=&
 \llbracket\phi\rrbracket\Rightarrow\llbracket\psi\rrbracket,\\
\llbracket\forall y<t\,\phi\rrbracket&=&
 \operatorname{All}(t,\lambda y.\llbracket\phi\rrbracket),\\
\llbracket\exists y<t\,\phi\rrbracket&=&
 \operatorname{Any}(t,\lambda y.\llbracket\phi\rrbracket).
\end{array}
\]
Again, each apparent lambda is eliminated at compile time because the formula
is fixed.

\begin{theorem}[Bounded Booleanization]
\label{thm:booleanization}
For every fixed bounded formula \(\phi(\vec x)\), there is a primitive-
recursive Boolean symbol \(b_\phi\) such that iPRA proves
\[
      \phi(\vec x)\quad\leftrightarrow\quad b_\phi(\vec x)=1.
\]
The equation \(b_\phi(\vec x)=1\) is quantifier free and can be used as an
induction invariant.
\end{theorem}

\begin{proof}
Proceed by structural induction on \(\phi\).  Atomic predicates are represented
by their characteristic functions.  The propositional cases follow from the
finite truth tables, each of which is provable by numerical case analysis on
Boolean inputs.  For bounded universal and existential quantifiers, apply the
fold-reflection lemma to the characteristic function supplied by the
induction hypothesis.  Primitive recursiveness follows by closure under
composition and primitive recursion.  The final equation is atomic.
\end{proof}

\subsection{Intuitionistic implication}

The implication clause deserves comment.  For arbitrary propositions, a
Boolean truth value would classically decide the antecedent.  Here the
antecedent and consequent are bounded formulas with explicit characteristic
functions.  iPRA proves that each is equivalent to a Boolean equation; Boolean
equality is decidable by finite case analysis.  Therefore
\[
 (b_\phi=1\to b_\psi=1)
 \leftrightarrow
 (b_\phi\Rightarrow b_\psi)=1
\]
is intuitionistically derivable.  No global law of excluded middle is used.

\subsection{Why arbitrary formula codes are excluded}

A uniform function taking an arbitrary code of a bounded formula and
returning its characteristic function would be an evaluator for a growing
program language.  A clocked syntactic evaluator is possible, but a total PR
evaluator for all PR programs is not.  The framework therefore uses one of two
modes:
\begin{enumerate}
\item compile a fixed formula at the Lean meta-level into a fixed \(\PR\)
symbol and proof; or
\item evaluate a coded formula for a supplied bounded clock and verify its
trace.
\end{enumerate}
This restriction is not a loss for ordinary theorem formalization, where the
formulas are known in advance.

\subsection{Lean formalization path}

The supplied file \code{BoundedFormulaCompiler.lean} defines a meta-syntax of
fixed bounded formulas and records the intended compilation functions.  A
complete version should proceed in four layers:
\begin{enumerate}
\item add canonical Boolean symbols and prove their equations;
\item define \code{allFold}, \code{anyFold}, and \code{leastFold} as
\(\PR\) terms;
\item prove their object-language reflection theorems by \code{PRA.qfInd};
\item recursively produce a pair consisting of \(b_\phi\) and a
\(\Prf\,\PRAAx\) proof of the equivalence.
\end{enumerate}
The fold layer is the only place requiring induction.  The structural compiler
then uses ordinary propositional proof constructors and substitution.

\section{Fix precision first, then quantify}
\label{sec:fixedprecision}

\subsection{The quantifier-order problem}

Many packets describe a construction that must work at every requested
precision \(p\).  A tempting induction predicate is
\[
   \forall p\;\Inv(i,p,u).
\]
This is not quantifier free, so it cannot be fed to \code{PRA.qfInd}.  Even when
one can encode the universal quantifier by a bounded fold, the bound may depend
on the loop counter in an inconvenient way.  The correct proof order is:
\begin{enumerate}
\item take a single arbitrary precision \(p\) as a free parameter;
\item compute all finite queries needed at that precision;
\item perform quantifier-free induction with \(p\) free;
\item universally introduce \(p\) after the induction is complete.
\end{enumerate}

\subsection{Finite query tables}

Suppose a construction at output precision \(p\) makes finitely many oracle
queries to its input presentation.  Let
\[
   Q(p)=\langle q_0(p),\ldots,q_{m(p)-1}(p)\rangle
\]
be the primitive-recursive query table, and let
\(A(p)\) encode the corresponding finite answers.  The condition that all
answers are valid is the Boolean fold
\[
   \operatorname{TableOK}(x,p,A)
   =\operatorname{All}\bigl(m(p),
       j\mapsto \operatorname{AnswerOK}(x,q_j(p),A_j)\bigr).
\]
It is one quantifier-free equation.  The loop invariant may mention \(p,Q(p),A\)
freely without quantifying over other precisions.

\begin{theorem}[Fixed-precision induction]
Let \(\Phi(i,p,u)\) be a fixed bounded formula and let
\(b_\Phi(i,p,u)\) be its Booleanization.  If iPRA proves the base and successor
steps uniformly with free parameters \(p,u\), then iPRA proves
\[
      \forall p\,\forall u\,\forall i\;\Phi(i,p,u).
\]
The only induction used is quantifier-free induction on the equation
\(b_\Phi(i,p,u)=1\).
\end{theorem}

\begin{proof}
Work with \(p,u\) as free variables.  By \cref{thm:booleanization}, replace
\(\Phi\) by the equivalent equation \(b_\Phi=1\).  Apply quantifier-free
induction on \(i\), using the supplied base and successor derivations after
transport across the Booleanization equivalence.  This yields
\(\forall i\,b_\Phi(i,p,u)=1\) with \(p,u\) still free.  Convert back to
\(\Phi\), then use universal introduction first for \(u\) and then for \(p\).
The eigenvariable conditions hold because the context was shifted exactly as
required by the checked \code{Prf.allI} constructor.
\end{proof}

\subsection{Simultaneous finite windows}

Sometimes a theorem needs a whole window of precisions
\(p,p+1,\ldots,p+k\), for example to verify a tail estimate or a finite
coherence condition.  Do not introduce \(k+1\) independent quantified
invariants.  Extend the query table to contain the union of all queries needed
by the window and fold their validity into one Boolean.  The window length may
be a free parameter if its loop is bounded and compiled into the Boolean
state.

\subsection{Changing parameters in recursive algorithms}

An algorithm may refine a precision or alter a denominator cap at each step.
Let \(d(i,p,u)\) be the demand at stage \(i\), computed by primitive recursion.
The invariant is not ``all demands have been met'' as a quantified formula; it
is a Boolean statement about the finite state containing the answers already
collected for the current demand table.  A reverse-demand pass, described in
\cref{sec:dag}, computes a sufficiently strong initial table so that no new
unbounded oracle call is needed during the forward run.

\subsection{Examples}

\paragraph{Taylor approximation.}
Fix a requested error \(2^{-p}\).  Compute a derivative order \(N(p)\), a
radius margin, and all coefficient precisions before proving the finite Taylor
identity and remainder bound by induction on \(N(p)\).  Quantify over \(p\)
last.

\paragraph{Root clusters.}
Fix a disk radius and coefficient precision.  The finite root-count verifier
uses one grid and one finite boundary table.  It does not induct on a predicate
saying that the count is correct at every smaller radius.

\paragraph{Simplicial approximation.}
Fix the output subdivision level.  Compute the source subdivision and carrier
cloud required by the modulus.  Prove the chain-map and contiguity identities
for those finite complexes, then quantify over subdivision level.

\subsection{Why this resolves a recurring packet difficulty}

The difficulty reports often identify an induction that appears to need a
quantifier over all precisions.  In almost every such case, the quantifier is
an artifact of proving too strong a statement inside the loop.  The finite
computation only uses finitely many observations for a fixed requested
output.  Encoding exactly that finite-use table restores a quantifier-free
invariant and produces a theorem that is extensionally just as strong after
universal introduction.

\section{The proof-producing finite-trace compiler}
\label{sec:tracecompiler}

\subsection{Trace systems}

A finite computation is internalized by fixed primitive-recursive symbols.
For clarity, consider one input parameter \(u\); tuples of parameters are
paired.

\begin{definition}[Trace system]
A trace system consists of fixed symbols
\[
\begin{array}{rclcrcl}
 \Good&:&\N\to\Bool, && \operatorname{init}&:&\N\to\N,\\
 \operatorname{bound}&:&\N\to\N, &&
 \operatorname{step}&:&\N^3\to\N,\\
 \Inv&:&\N^3\to\Bool, &&
 \Accept&:&\N^2\to\Bool.
\end{array}
\]
The run is the primitive recursion
\begin{align*}
 \Run(0,u)&=\operatorname{init}(u),\\
 \Run(Si,u)&=\operatorname{step}(i,\Run(i,u),u).
\end{align*}
\end{definition}

All predicates are represented by equations to the Boolean value one.  We use
the abbreviations
\[
\begin{aligned}
 G(u)&:\equiv \Good(u)=1,\\
 I(i,s,u)&:\equiv \Inv(i,s,u)=1,\\
 A(s,u)&:\equiv \Accept(s,u)=1.
\end{aligned}
\]

\subsection{Local proof interface}

The compiler accepts three closed object-language derivations:
\begin{align}
 \mathsf{base}:&\quad []\vdash[\PRAAx]
       G(u)\to I(0,\operatorname{init}(u),u),\label{eq:tracebase}\\
 \mathsf{step}:&\quad []\vdash[\PRAAx]
       G(u)\to I(i,s,u)\to
       I(Si,\operatorname{step}(i,s,u),u),\label{eq:tracestep}\\
 \mathsf{final}:&\quad []\vdash[\PRAAx]
       G(u)\to I(\operatorname{bound}(u),s,u)\to A(s,u).
       \label{eq:tracefinal}
\end{align}
Their displayed free variables are represented by de Bruijn indices.  The
conclusion to construct is
\begin{equation}
 []\vdash[\PRAAx]
 \forall u\,[G(u)\to
 A(\Run(\operatorname{bound}(u),u),u)].
 \label{eq:traceresult}
\end{equation}

\begin{theorem}[Finite-trace reflection]
\label{thm:tracereflection}
For every fixed trace system, the three derivations
\eqref{eq:tracebase}--\eqref{eq:tracefinal} primitive-recursively determine a
derivation of \eqref{eq:traceresult}.
\end{theorem}

\subsection{Human proof of trace reflection}

We give the proof at the same granularity as the intended Lean constructor.
No semantic induction outside iPRA is used.

\paragraph{Step 1: open the input and assume the guard.}
Work in one free variable \(u\) and context \([G(u)]\).  Because
\(\mathsf{base}\) is closed, substitute the term \(u\) for its free variable
and weaken it to this context.  Apply implication elimination to obtain
\[
    [G(u)]\vdash I(0,\operatorname{init}(u),u).
\]
The primitive-recursion zero equation gives
\[
    [G(u)]\vdash
    \Run(0,u)=\operatorname{init}(u).
\]
Use Leibniz transport in the state position of \(I(0,-,u)\) to derive
\begin{equation}
   [G(u)]\vdash I(0,\Run(0,u),u).
   \label{eq:invbase}
\end{equation}

\paragraph{Step 2: construct the induction step.}
The induction formula, with free variables \((i,u)\), is
\[
      \Theta(i,u):\equiv I(i,\Run(i,u),u).
\]
It is an equation to one and hence quantifier free.  For the successor step,
work in the shifted context
\([\Theta(i,u),G(u)]\).  Instantiate \(\mathsf{step}\) with the terms
\[
     i,\\ \Run(i,u),\\ u.
\]
Eliminate the two implications using the guard and the induction hypothesis.
This yields
\[
 I(Si,\operatorname{step}(i,\Run(i,u),u),u).
\]
The primitive-recursion successor equation gives
\[
 \Run(Si,u)=\operatorname{step}(i,\Run(i,u),u).
\]
Transport in the state argument, reversing the equality when necessary with
\code{Prf.eqSymm}, to obtain
\begin{equation}
 [\Theta(i,u),G(u)]\vdash\Theta(Si,u).
 \label{eq:invstep}
\end{equation}

\paragraph{Step 3: invoke quantifier-free induction.}
Apply \code{PRA.qfInd} to \(\Theta\), using \eqref{eq:invbase} and
\eqref{eq:invstep}.  The result is
\begin{equation}
       [G(u)]\vdash \forall i\,I(i,\Run(i,u),u).
       \label{eq:allinv}
\end{equation}
The context shifting required by the induction rule is exactly the shifting
already built into \code{Prf.allI} and \code{PRA.qfInd}; no informal
renaming convention is needed.

\paragraph{Step 4: specialize at the computed bound.}
Use universal elimination in \eqref{eq:allinv} with the term
\(\operatorname{bound}(u)\).  We obtain
\begin{equation}
 [G(u)]\vdash
 I(\operatorname{bound}(u),
   \Run(\operatorname{bound}(u),u),u).
 \label{eq:boundinv}
\end{equation}

\paragraph{Step 5: apply the final contract.}
Instantiate \(\mathsf{final}\) with
\[
 s=\Run(\operatorname{bound}(u),u)
 \quad\text{and}\quad u=u.
\]
Weaken to the guard context and eliminate the two implications using the guard
and \eqref{eq:boundinv}.  This gives
\[
  [G(u)]\vdash
  A(\Run(\operatorname{bound}(u),u),u).
\]
Discharge \(G(u)\) by implication introduction and then close \(u\) by
universal introduction.  This is exactly \eqref{eq:traceresult}.

\subsection{Why the construction is primitive recursive on proofs}

The derivation is assembled from a fixed finite pattern of constructors plus
the three input derivations.  It performs no search through proofs.  Term and
formula substitutions are structural recursive functions already defined in
the checked library.  Hence, at the meta-level, \code{traceReflection} is an
ordinary total Lean definition.

\subsection{Multiple state components and certificate checks}

A single natural state can encode a tuple containing:
\begin{itemize}
\item program variables;
\item a current query-table index;
\item accumulated Boolean validity bits;
\item an error numerator and denominator;
\item a lexicographic descent measure;
\item a bounded certificate cursor.
\end{itemize}
The invariant remains a single Boolean equation checking all components.  A
large conjunction should be compiled as a Boolean fold, not represented by a
formula with a variable number of logical connectives.

\subsection{Nested loops}

Nested bounded loops may be flattened into one trace by pairing loop counters,
or handled compositionally.  Flattening is useful when one global invariant is
simpler than a hierarchy of intermediate theorems.  Composition is useful when
local loop proofs are reusable.  Both routes stay inside quantifier-free
induction because every loop is bounded and each local invariant has been
Booleanized.

\subsection{Termination by a ranked measure}

Some algorithms are described as ``repeat until done'' but come with a
primitive-recursive rank \(r(s)\) that strictly decreases when the terminal
predicate fails.  If an initial bound \(R(u)\) on the rank is supplied, compile
the algorithm to exactly \(R(u)\) steps, making terminal states stutter.  The
invariant records
\[
 \operatorname{done}(s_i)=1
 \quad\text{or}\quad r(s_{i+1})<r(s_i).
\]
A bounded induction proves that a nonterminal state cannot persist beyond
\(R(u)\) decreases.  Thus ranked descent is a special case of trace reflection,
not an additional induction principle.

\subsection{Lean realization}

The file \code{ObjectTraceCompiler.lean} contains arity-correct builders
\code{app1}, \code{app2}, and \code{app3}; the trace-system structure; all five
formula builders; the quantifier-free lemma; the run equations; and exact
signatures for four milestones:
\begin{enumerate}
\item \code{trace_base};
\item \code{trace_step};
\item \code{trace_all};
\item \code{trace_final}.
\end{enumerate}
The main definition \code{traceReflection} is then a short composition of
those milestones.  The remaining \code{sorry} obligations are substitution
normalizations and equality transports, not missing mathematical ideas.  The
last theorem
\[
 \code{traceReflectionNative}
\]
is already one application of the checked \code{pra_iff_native} bridge once
\code{traceReflection} is complete.

\subsection{Consequence for the proof packets}

Any packet whose main proof has been reduced to a fixed bounded algorithm with
local iPRA derivations can use \cref{thm:tracereflection}.  It no longer needs a
custom meta-argument that the algorithm ``ought to be formalizable.''  The
packet must provide exactly three kinds of evidence:
\begin{enumerate}
\item primitive-recursive symbols for the state transition and bound;
\item a quantifier-free invariant strong enough for the theorem;
\item closed object-language derivations of base, preservation, and final
soundness.
\end{enumerate}
This is the smallest reusable bridge between finite mathematics and the
checked logical backend.

\section{Contract DAGs and reverse demand propagation}
\label{sec:dag}

\subsection{Why local composition is not enough}

A large formalization is rarely a single loop.  It is a directed acyclic graph
of constructions: coefficient approximation feeds root isolation and an error
estimate; both feed a contour certificate; several branches share the same
input observations.  Naive nested composition recomputes demands and may hide
an illegal adaptive oracle call.  A contract-DAG compiler separates two
passes:
\begin{enumerate}
\item a reverse pass determines every observation needed for the requested
output;
\item a forward pass executes each node once and records one bounded trace.
\end{enumerate}

\subsection{DAG model}

Let nodes be topologically ordered \(0,\ldots,N-1\).  Node \(v\) has parents
\(p_0(v),\ldots,p_{k_v-1}(v)<v\).  Its contract contains:
\begin{itemize}
\item an output query domain \(Q_v\);
\item for each output query \(q\), a parent-demand tuple
\(D_v(q)=(q_0,\ldots,q_{k_v-1})\);
\item a primitive-recursive evaluator \(E_v\);
\item a local validity predicate \(C_v\);
\item a local \(\Prf\,\PRAAx\) derivation of contract preservation.
\end{itemize}
If query domains are numerical precision levels, join is maximum.  For richer
finite queries, assume a primitive-recursive join \(q\sqcup q'\) with proofs
that it dominates both.  The displayed \(D_v\) is static: it depends on the
output query and fixed finite metadata, not on an as-yet uncomputed approximate
value.  Data-dependent algorithms must provide a certified static majorant, a
finite preliminary summary stage, or the joined demands of all unresolved
branches.  This condition is what makes the reverse pass genuinely prior to
the forward pass.

\subsection{Reverse demand algorithm}

Initialize every node demand to the least query \(\bot\), except the output
nodes, which receive the user's requested queries.  Scan nodes in reverse
topological order.  At node \(v\), compute \(D_v(q_v)\) and join each component
into the demand of the corresponding parent.  Because the graph is finite and
the scan length is \(N\), this is primitive recursive.

\begin{theorem}[Soundness of reverse demand]
\label{thm:demandsound}
Let \(q_v\) be the demands produced by the reverse pass.  If each source
presentation supplies a valid observation at its computed demand, then every
node \(v\) can be evaluated at demand \(q_v\), and every requested output
contract holds.
\end{theorem}

\begin{proof}
Induct on the topological order.  Source nodes satisfy their contracts by the
assumption.  For a non-source node \(v\), the reverse pass joined the demand
\(D_v(q_v)_j\) into each parent \(p_j(v)\).  Hence the observation produced for
that parent is at least as strong as the node's local requirement.  Restrict it
if necessary, apply the local contract derivation, and obtain a valid output at
\(q_v\).  Continue through the finite order.  The requested outputs have the
queries with which the reverse pass was initialized, so their contracts hold.
\end{proof}

\subsection{Fan-out and sharing}

Suppose a node feeds two consumers requiring precisions \(p\) and \(q\).  The
producer is evaluated once at \(p\sqcup q\).  Restriction maps supply the two
weaker observations.  This is more than an optimization: it ensures that both
consumers see coherent data from one finite observation rather than unrelated
witnesses.

For relational presentations, sharing does not force a single exact candidate.
The producer may return a cloud at the joined demand, and each consumer uses a
compatible restriction cloud.  A global viability pass selects a finite
compatible path only when needed.

\subsection{The global forward trace}

Encode in the trace state:
\begin{itemize}
\item the current node index;
\item the complete demand table;
\item the finite table of already computed observations;
\item one accumulated validity Boolean;
\item optional local certificate cursors.
\end{itemize}
At each step, the transition evaluates the next node from its parent entries
and updates the validity bit by conjunction with \(C_v\).  The invariant says:
\begin{enumerate}
\item every completed node has a well-formed observation at its assigned
demand;
\item every completed local contract is true;
\item untouched table entries retain the default value;
\item the demand table equals the result of the reverse pass.
\end{enumerate}
Each clause is a bounded Boolean fold.  The invariant is one equation and can
be fed to \cref{thm:tracereflection}.

\begin{theorem}[Certified DAG compilation]
A finite acyclic graph of proof-carrying contracted transducers with static or
certifiably majorized demands compiles to one trace system.  From the local
\(\Prf\,\PRAAx\) contract proofs, iPRA derives the global output contract at
every requested finite query.
\end{theorem}

\begin{proof}
The reverse-demand correctness theorem supplies the semantic invariant.
Booleanize its bounded table clauses.  The base proof verifies the initialized
tables.  The step proof is a finite case distinction on the current node,
followed by the corresponding local contract proof and table-update equations.
The final proof projects the accumulated output entries and their validity
bits.  Apply finite-trace reflection and universally quantify the requested
queries and source inputs.
\end{proof}

\subsection{Guarded branches}

A branch node may be compiled in three safe ways.
\begin{enumerate}
\item If the condition is a decidable finite Boolean, execute exactly one
branch.
\item If an approximate quantity is separated from the branch boundary by a
positive margin, refine until the finite guard certifies one side, then execute
that branch.
\item Without separation, execute or describe both finite branches and return
a relational union whose observable output is proved stable.
\end{enumerate}
An exact branch on the sign of an arbitrary real name is not permitted.

\subsection{Loops in a program graph}

A bounded loop is represented as one trace node whose local proof is itself
obtained from \cref{thm:tracereflection}.  Alternatively, the compiler can
flatten all loops into the global state.  The modular route creates smaller
local proof obligations; the flattened route permits stronger sharing.  Both
are conservative because the proof terms are composed by substitution and cut.

\subsection{Lean outline}

The supplied \code{CertifiedProgram.lean} sketches a syntax with assignments,
checks, sequential composition, bounded loops, and guarded or relational
branches.  Completion should proceed by:
\begin{enumerate}
\item replacing abbreviated semantic fields by formula builders and
\(\Prf\,\PRAAx\) derivations;
\item defining a topological list representation with a bounded well-formedness
checker;
\item implementing the reverse pass as a \(\PR\) fold;
\item proving the demand invariant by quantifier-free induction;
\item compiling the forward table update to a \code{PRTraceSystem};
\item invoking \code{traceReflection} once.
\end{enumerate}
No higher-order evaluator is needed: each certified program is a fixed Lean
value from which fixed \(\PR\) symbols are constructed.

\section{Finite relational witnesses and backward viability}
\label{sec:relational}

\subsection{The coherence problem}

For each precision \(p\), a proof may produce at least one valid finite witness,
but witnesses at different precisions may be incompatible.  Choosing the least
candidate independently at each level does not solve the problem.  It may jump
between branches and fail to define a name.  Conversely, appealing to an
infinite path through a finitely branching tree would import a choice or
compactness principle stronger than the finite theorem being proved.

The remedy is to select only a finite path to the finite horizon currently
requested.

\subsection{Finite tower}

Fix a horizon \(h\).  For each level \(p\le h\), let
\[
 C_p=\{a<B_p:\Cand(p,a)=1\}
\]
be a finite candidate cloud.  Let \(E_p(a,b)\) be the compatibility relation
from level \(p\) to \(p+1\).

Define backward viability relative to \(h\):
\begin{align*}
 V_h(h,a)&:\equiv \Cand(h,a),\\
 V_h(p,a)&:\equiv \Cand(p,a)\wedge
   \exists b<B_{p+1}\,[E_p(a,b)\wedge V_h(p+1,b)].
\end{align*}
All existential quantifiers are bounded and are implemented by Boolean folds.

\begin{theorem}[Viability/path equivalence]
\label{thm:viability}
For \(p\le h\), \(V_h(p,a)\) holds if and only if there is a compatible path
\[
 a=a_p,a_{p+1},\ldots,a_h
\]
with \(a_j\in C_j\) and \(E_j(a_j,a_{j+1})\) for every \(j<h\).
\end{theorem}

\begin{proof}
Use backward induction on \(h-p\).  At \(p=h\), viability is exactly candidate
validity and the one-element path suffices.  For the induction step, unfold
\(V_h(p,a)\).  A viable successor \(b\) exists; by the induction hypothesis it
extends to a compatible tail, which we prepend by \(a\).  Conversely, the
second element of any compatible full path is a viable successor by the
induction hypothesis, so the defining bounded existential holds.
\end{proof}

\subsection{Least viable path}

Define \(a_p\) to be the least viable candidate at the initial level.  Given
\(a_j\), choose the least \(a_{j+1}<B_{j+1}\) satisfying
\[
 E_j(a_j,a_{j+1})\wedge V_h(j+1,a_{j+1}).
\]
By \cref{thm:viability}, such a candidate exists whenever the initial
viability test succeeds.  Bounded least search is primitive recursive.

\begin{corollary}[Finite dependent choice]
For every finite horizon, a nonempty viable tower has a primitive-recursively
selected compatible path.  This construction uses only bounded search and
quantifier-free induction; it does not imply an infinite choice principle.
\end{corollary}

\subsection{Stable quotient erasure}

Often the theorem does not need the internal candidate; it needs an observable
\(W_h(a_0,\ldots,a_h)\).  If every viable path has the same observable, then
choice is irrelevant.

\begin{theorem}[Exact stable erasure]
Assume that for any two viable paths \(\alpha,\beta\) to horizon \(h\),
\[
       W_h(\alpha)=W_h(\beta).
\]
Then the output computed from the least viable path is independent of every
tie-breaking convention and is extensionally characterized by the finite
relation alone.
\end{theorem}

\begin{proof}
The least viable path is viable.  Apply the stability assumption to it and to
an arbitrary viable path.  Their observables are equal.  Since all quantifiers
range over bounded path codes, the argument Booleanizes into a finite proof.
\end{proof}

Exact equality can be replaced by a metric diameter bound:
\[
 d(W_h(\alpha),W_h(\beta))\le\varepsilon_h.
\]
Then every tie-breaking convention yields outputs within \(\varepsilon_h\).
This is the natural statement for approximate roots, eigenvalues, local
solutions, and geometric selections.

\subsection{Changing horizons}

The path selected for horizon \(h+1\) need not extend the path selected for
horizon \(h\).  This does not matter if the level-\(p\) candidates lying on
full viable paths to either horizon have small diameter.  Define the selected
observable at precision \(p\) using a horizon \(H(p)\ge p\).  Suppose:
\begin{enumerate}
\item full-horizon viable candidates at level \(p\) have diameter at most
\(\delta_p\);
\item every compatibility edge changes the represented value by at most
\(\mu_p\);
\item a primitive-recursive tail bound controls
\(\sum_{j\ge p}\mu_j\).
\end{enumerate}
Then selections made with different future horizons are close.  The next
section turns this fact into a regular name.

\subsection{Applications}

\paragraph{Polynomial roots.}
At precision \(p\), candidates are finite rational disks or root clusters.
Edges express refinement and preservation of multiplicity.  At collisions,
there is no ordered label; symmetric observables or cluster data are stable.
A separation certificate permits exact labels locally.

\paragraph{Local charts.}
Candidates are chart indices and approximate coordinates.  Edges assert that
overlap transition maps agree within a rational budget.  A downstream theorem
uses a stable observable, not a globally canonical chart.

\paragraph{Finite sheaf data.}
Candidates are local sections satisfying equations on a finite cover.  Edges
are restriction compatibility.  A gluing certificate may return a cloud of
global representatives; quotient equality erases nonuniqueness.

\paragraph{Mordell--Weil descent.}
Candidates can encode rational points on covering spaces together with maps to
the base curve.  The framework verifies compatibility and descent.  It does
not manufacture rational representatives from abstract Selmer classes; that
remains a producer obligation.

\subsection{Lean plan}

The file \code{RelationalWitnesses.lean} defines finite bounds, validity,
edges, and viability at the semantic level and records object-language targets.
A complete migration should:
\begin{enumerate}
\item encode the level/horizon dynamic program as fixed \(\PR\) symbols;
\item prove the viability recursion equations;
\item Booleanize path validity;
\item prove both directions of \cref{thm:viability} by quantifier-free
induction on the remaining horizon;
\item define least viable successor with the shared bounded-search library;
\item formulate stability as a local finite theorem and compose it with the
selected path proof.
\end{enumerate}

\section{Regular names, shrinking viability, and guarded operations}
\label{sec:namesguards}

\subsection{Regular rational names}

A convenient constructive real presentation is a rational sequence
\(q:\N\to\Q\) with a fixed regularity law, for example
\[
       |q_{p+1}-q_p|\le 2^{-p}.
\]
The law implies a primitive-recursive tail estimate
\[
       |q_m-q_p|\le 2^{1-p}
       \qquad(m\ge p).
\]
All statements are finite rational inequalities.  A completed real is never
needed as a set-theoretic object inside iPRA.

More generally, a name may come with a modulus \(M(p)\) such that
\[
 m,n\ge M(p)\quad\Longrightarrow\quad |q_m-q_n|\le 2^{-p}.
\]
The modulus must be explicit.  Mere convergence or metastability is not enough
for operations requiring a uniform tail.

\subsection{Names from relational towers}

Let \(v(p,a)\in\Q\) be the value of candidate \(a\) at level \(p\).  For an
increasing primitive-recursive horizon \(H(p)\), choose the least candidate at
level \(p\) that lies on a viable path to \(H(p)\); call its value \(q_p\).
Assume:
\begin{align}
 \diam\{v(p,a):V_{H}(p,a)\}&\le\delta_p
   &&\text{for every }H\ge H(p),\label{eq:diambound}\\
 E_p(a,b)&\Longrightarrow |v(p,a)-v(p+1,b)|\le\mu_p,
   \label{eq:motionbound}
\end{align}
and there is a primitive-recursive tail bound
\[
       \sum_{j=p}^{m-1}\mu_j\le T(p)
       \quad(m>p),
       \qquad T(p)\to0
\]
in an explicit numerical form.

\begin{theorem}[Shrinking viability gives a name]
\label{thm:shrinkingname}
For \(m>p\),
\[
      |q_m-q_p|\le \delta_p+T(p)+\delta_m.
\]
Consequently, if the right-hand side is at most \(2^{-r}\) whenever
\(p,m\ge M(r)\), then \((q_p)\) is a Cauchy name with modulus \(M\).
\end{theorem}

\begin{proof}
Take a viable path witnessing the choice at horizon \(H(m)\).  Its level-\(p\)
candidate, call it \(a_p\), remains viable after truncating that path at
horizon \(H(p)\), because \(H(m)\ge H(p)\).  The candidate defining \(q_p\) is
also viable to \(H(p)\).  Apply \eqref{eq:diambound} at the common horizon
\(H(p)\) to conclude that \(q_p\) and \(v(p,a_p)\) differ by at most
\(\delta_p\).  Follow the longer path from level \(p\) to level \(m\);
telescoping \eqref{eq:motionbound} bounds the change by \(T(p)\).  At level
\(m\), the candidate on this path and the independently selected candidate
defining \(q_m\) are both viable to the common horizon \(H(m)\), so their
values differ by at most \(\delta_m\).  Add the three rational inequalities.
Every path involved is finite, so the proof is a bounded certificate argument.
\end{proof}

This theorem avoids selecting one infinite branch.  Each \(q_p\) uses a finite
horizon; shrinking diameter makes the horizon dependence observationally
irrelevant.

\subsection{Three-way guards}

Exact binary branching on approximate data is generally discontinuous.  Given
a rational approximation \(q\) to a quantity \(x\), an error bound
\(|x-q|\le e\), and a positive margin \(\eta\), define
\[
 \operatorname{guard}(q,e,\eta)=
 \begin{cases}
  -,&q+e< -\eta,\\
  +,&q-e> \eta,\\
  ?,&\text{otherwise}.
 \end{cases}
\]
The output is a finite tag computed by rational comparison.

\begin{lemma}[Guard soundness]
If the guard returns \(+\), then \(x>\eta\).  If it returns \(-\), then
\(x<-\eta\).  If it returns \(?\), then the approximation has not certified a
side; no conclusion is forced.
\end{lemma}

\begin{proof}
For the positive case, \(x\ge q-e>\eta\).  The negative case is symmetric.
The unknown case is deliberately information-preserving.
\end{proof}

As precision increases, the error budget decreases.  An unknown guard may
later resolve, but a certified positive or negative guard remains sound.  When
the theorem assumes an apartness margin \(|x|>\eta\), sufficiently strong
precision guarantees resolution.

\subsection{Guarded reciprocal}

Let a regular name \(q_p\) represent \(x\), with
\(|x-q_p|\le e_p\).  Suppose an explicit rational \(\eta>0\) satisfies
\(|x|\ge\eta\).  Choose \(p\) with \(e_p\le\eta/2\).  Then
\(|q_p|\ge\eta/2\), so rational division by \(q_p\) is safe.

For two approximations \(q,r\) with
\(|x-q|,|x-r|\le e\le\eta/2\),
\[
 \left|\frac1q-\frac1r\right|
 =\frac{|q-r|}{|qr|}
 \le \frac{2e}{(\eta/2)^2}
 =\frac{8e}{\eta^2}.
\]
Thus, to obtain reciprocal precision \(2^{-k}\), it suffices to request an
input error at most \(2^{-k}\eta^2/8\).  This yields a primitive-recursive
backward demand once \(\eta\) is coded rationally.

\begin{theorem}[Guarded reciprocal]
There is a contracted transducer taking a regular rational name for \(x\), a
positive rational \(\eta\), and a certificate \(|x|\ge\eta\), and returning a
regular name for \(1/x\).  Its demand transformer is primitive recursive and
its correctness proof uses only rational inequalities.
\end{theorem}

The same pattern handles matrix inversion from a determinant or singular-value
margin, Gaussian pivots from a pivot margin, chart selection from a distance-to-
boundary margin, logarithms on positive intervals, and isolated algebraic
branches.

\subsection{Root and eigenvalue labels}

At a collision, an ordered list of roots or eigenvectors is not continuous.
The correct unconditional outputs are:
\begin{itemize}
\item a multiset or cluster cloud;
\item symmetric functions;
\item invariant subspaces or spectral projectors for separated clusters;
\item exact labels only with a separation gap.
\end{itemize}
If disks around roots are pairwise separated by \(\eta\), the least-disk
ordering is stable under perturbations smaller than \(\eta/3\).  The gap is a
guard and the label-selection proof is finite.  Without it, the framework
retains the relational cloud.

\subsection{Metastability versus a modulus}

A metastability statement has the form
\[
  \forall k\,\forall g\,\exists N\le B(k,g)\,
  \forall m,n\in[N,N+g(N)]\quad |q_m-q_n|\le2^{-k}.
\]
For each fixed \(k,g\), this is a finite-search theorem and fits the framework.
It does not supply a tail modulus valid for every \(m,n\ge N\).  Therefore it
supports finite-window conclusions but not, by itself, a regular name or
operations such as guarded reciprocal.  The type of the interface must retain
this distinction.

\section{Effective compactness and local-to-global constructions}
\label{sec:compact}

\subsection{Finite nets in place of abstract compactness}

Let \((X,d)\) be represented at finite precision.  An \(\varepsilon\)-net is a
finite list \(x_0,\ldots,x_{N-1}\) together with a coverage certificate
\[
       \forall x\in X\;\exists j<N\quad d(x,x_j)\le\varepsilon.
\]
In iPRA the displayed quantifiers are not taken over an abstract completed
space.  The presentation supplies a finite observation relation strong enough
to verify the coverage statement needed by the theorem.  For a rational box,
for example, a grid and coordinate rounding algorithm provide the certificate.
For a compact metric presentation, the net producer is part of the input
interface.

\subsection{Approximate minimization}

Suppose \(f:X\to\Q\) has a modulus
\[
 d(x,y)\le\varepsilon\quad\Longrightarrow\quad
 |f(x)-f(y)|\le\omega(\varepsilon).
\]
Evaluate \(f\) on the finite net and choose the least index \(j_*\) attaining
the minimum rational value.

\begin{theorem}[Finite-net minimum transfer]
For every represented \(x\in X\),
\[
       f(x_{j_*})\le f(x)+2\omega(\varepsilon)+2e,
\]
where \(e\) bounds each finite evaluation error.  In particular, the selected
net point is an explicit approximate minimizer.
\end{theorem}

\begin{proof}
Choose a net point \(x_j\) within \(\varepsilon\) of \(x\).  Let
\(\widetilde f_j\) denote the computed rational approximation, with
\(|\widetilde f_j-f(x_j)|\le e\).  By finite minimality,
\(\widetilde f_{j_*}\le\widetilde f_j\).  Hence
\[
 f(x_{j_*})
 \le \widetilde f_{j_*}+e
 \le \widetilde f_j+e
 \le f(x_j)+2e
 \le f(x)+\omega(\varepsilon)+2e.
\]
The stated slightly looser bound allows asymmetric input/output moduli and is
closed under the usual error-budget composition.  Every choice is bounded
least search.
\end{proof}

\subsection{Finite covers}

An effective finite cover consists of:
\begin{itemize}
\item finitely many chart or open-set codes \(U_i\);
\item a Boolean membership or positive-margin predicate at the required
precision;
\item a coverage certificate for the finite observation domain;
\item transition data on pairwise overlaps;
\item a Lebesgue-number or refinement margin when exact chart selection is
needed.
\end{itemize}
The framework never extracts a finite subcover from an arbitrary open cover.
It consumes or constructs a finite cover from explicit quantitative data.

\subsection{Finite gluing}

Let local data \(s_i\) be given on a finite cover.  Suppose compatibility on
overlaps is a bounded decidable relation
\[
       R_{ij}(s_i,s_j)=1.
\]
A gluing theorem can have one of three outputs:
\begin{enumerate}
\item an exact global code, when a bounded algebraic constructor exists;
\item a finite cloud of global codes, with quotient-stable observable;
\item a regular name obtained from compatible approximate gluings and a
shrinking diameter bound.
\end{enumerate}

For affine algebraic data, gluing often reduces to a finite system of ring
equalities.  For simplicial data, it reduces to finite carrier and contiguity
relations.  For analytic local solutions, it reduces to rational coefficient
bounds and overlap estimates.

\begin{theorem}[Finite gluing by bounded search]
Assume a primitive-recursive bound \(B\) on global candidate codes and a
bounded verifier \(G(c,s_0,\ldots,s_{r-1})\) that is sound and complete for the
finite gluing problem.  Then the finite list of all global gluings is primitive
recursive.  If a stable observable is constant on that list, it is a
well-defined deterministic output.
\end{theorem}

\begin{proof}
Enumerate \(c<B\) and retain exactly those passing \(G\).  Soundness says every
retained code is a gluing; completeness says every gluing appears.  The list is
a bounded fold.  Constancy of the observable erases the choice of retained
code.  The complete list, rather than an unbounded existential, is the
constructive global object.
\end{proof}

\subsection{Local differential equations}

For a rational polynomial ODE on a rational box, a finite local solution may
be represented by a truncated power series plus a remainder bound.  A
contracted local solver consumes:
\begin{itemize}
\item coefficient bounds;
\item a Lipschitz or majorant constant;
\item a rational step size satisfying the contraction inequality;
\item an initial-value approximation and error budget.
\end{itemize}
It returns a finite coefficient table and a verified residual bound.  A global
solution on a finite time interval is obtained by a bounded number of local
steps.  The state invariant accumulates the error budget; no abstract Picard
fixed point is needed in the object language.

\subsection{Local algebraic branches}

An equation \(F(x,y)=0\) with a certified nonzero partial derivative
\(|\partial_yF|\ge\eta\) admits a guarded local branch.  A rational Newton or
bisection procedure produces a finite regular name; the derivative margin
controls uniqueness and demand propagation.  At singular points the framework
returns a finite relation of branches or consumes exact algebraic splitting
data.  This distinction addresses the recurring Puiseux difficulty: exact
branch coefficients are safe for exact algebraic presentations, but not as a
uniform operation on arbitrary Cauchy coefficients.

\subsection{Observational quotients and extensionality}

Suppose two presentation codes \(x,x'\) represent the same object by
\(x\equiv_p x'\) at all precisions.  A contracted transducer is extensional if
for every output query \(q\), input equivalence at \(\Need_F(q)\) implies
output equivalence at \(q\).  Composition preserves extensionality by
substituting the demand inequalities.  Thus quotient reasoning becomes an
ordinary theorem about finite observations; no quotient type or choice of
canonical representatives is required inside iPRA.

\subsection{Local-to-global principle}

The general local-to-global pattern is now finite:
\begin{enumerate}
\item produce a finite cover or decomposition;
\item solve each local problem with explicit contracts;
\item verify pairwise compatibility by bounded formulas;
\item enumerate or construct finite global candidates;
\item erase nonuniqueness by exact stability or a diameter modulus;
\item compile the entire bounded process with one trace invariant.
\end{enumerate}
This pattern applies to finite sheaves, charted manifolds, local analytic
branches, and many descent calculations.  What it does not do is turn an
ineffective existence theorem for a cover, bound, or local solution into an
effective producer.

\section{Reusable domain templates}
\label{sec:domains}

The framework is useful only if domain experts can recognize where their
ordinary proof enters it.  This section gives conversion templates rather than
claiming that every deep producer theorem has already been proved.

\subsection{Finite algebra and linear algebra}

\paragraph{Objects.}
Coded commutative rings, modules, matrices, polynomials, finite presentations,
and bounded monomial tables.

\paragraph{Primitive-recursive core.}
Matrix arithmetic, determinants by Leibniz or fraction-free elimination,
characteristic polynomials, polynomial evaluation, truncated power series,
finite Gröbner/Macaulay calculations at a supplied degree, and verification of
factorizations or syzygies.

\paragraph{Typical contracts.}
A matrix inverse consumes a determinant/pivot margin or an exact adjugate
certificate.  A spectral projector consumes a gap separating the target
cluster.  Polynomial roots are returned as clusters unless a separation
certificate permits labels.

\paragraph{What becomes straightforward.}
Fixed-dimension identities and variable-dimension identities with bounded
folds are ideal trace-compiler targets.  Cayley--Hamilton, determinant
identities, finite module calculations, and exact certificate checking all
reduce to polynomial identities and bounded sums.

\subsection{Constructive analysis}

\paragraph{Objects.}
Regular rational names, interval or ball names, finite coefficient tables,
moduli of continuity, derivative bounds, and tail estimates.

\paragraph{Primitive-recursive core.}
Rational interval arithmetic, finite sums and products, Taylor polynomials,
quadrature rules, Euler--Maclaurin sums, bisection, and finite contour meshes.

\paragraph{Contracts.}
Every limiting theorem must expose a modulus or a finite error majorant.  A
partial operation such as reciprocal or logarithm consumes an apartness or
positivity margin.  An oscillatory or branch-sensitive construction returns a
relation until a gap resolves it.

\paragraph{Residual producer obligations.}
A classical existence proof of a convergence constant does not automatically
produce a primitive-recursive bound.  The framework can verify any proposed
constant and compose its consequences, but it cannot erase ineffectivity.

\subsection{Number theory and Diophantine geometry}

\paragraph{Objects.}
Integers, rationals, ideals in finite presentations, bounded-height points,
finite descent data, local solubility certificates, and exact recurrence
sequences.

\paragraph{Primitive-recursive core.}
Divisibility, gcd, modular arithmetic, continued fractions, finite lattice
search, bounded rational-point enumeration, verification of descent maps, and
finite prime sums.

\paragraph{Deep interface.}
A Diophantine finiteness theorem is split into:
\begin{enumerate}
\item a producer of an explicit height bound or complete list;
\item a bounded enumerator and verifier;
\item a proof that the list is complete below the bound.
\end{enumerate}
The last two are generic and iPRA-friendly.  The first contains the deep
mathematics.  For Mordell--Weil descent, an abstract finite Selmer group is not
enough: rational representatives on covering spaces, or a relational
certificate explaining their absence, must be supplied.

\subsection{Analytic number theory}

\paragraph{Objects.}
Finite prime tables, rational approximations to logarithms and exponentials,
truncated Dirichlet sums, zero-free-region certificates, and explicit error
budgets.

\paragraph{Primitive-recursive core.}
Sieve tables, divisor sums, partial summation, finite Euler products,
Möbius transforms, and verification of contour or kernel estimates.

\paragraph{Quantifier discipline.}
A useful theorem states the exact dependence of the cutoff on all numerical
parameters.  Phrases such as ``for sufficiently large \(x\), uniformly in
\(q\)'' must be expanded into a primitive-recursive bound with the quantifier
order fixed before compilation.

\paragraph{Residual obligations.}
Ineffective exceptional-zero arguments or unspecified constants remain
producer obligations.  The framework can still formalize conditional versions
and every finite verification step.

\subsection{Combinatorics}

Finite combinatorics is the most direct domain.  Graphs, set partitions,
finite posets, words, tableaux, and recurrences are bounded data.  Enumeration,
dynamic programming, involutions, sign-reversing pairings, and generating-
function coefficient identities compile to bounded traces.  Asymptotic claims
need explicit remainder estimates, but exact recurrences and congruences do
not.

For example, Bell-number recurrences and Touchard congruences are finite fold
theorems.  Dobi\'nski-type formulas require a truncation bound and a rational
tail estimate.  The interface cleanly separates the exact combinatorial core
from the analytic producer.

\subsection{Probability}

\paragraph{Objects.}
Finite rational probability laws, bounded random variables, rational moments,
finite product spaces, and explicit tail or quantization bounds.

\paragraph{Primitive-recursive core.}
Expectation is a finite weighted sum.  Independence is a table identity.
Convolution, characteristic-function samples, and finite couplings are
bounded calculations.

\paragraph{Approximation interface.}
A general distribution is represented by finite rational laws at precision
\(p\), together with a transport, tail, or anti-concentration modulus.  CDF
comparisons at discontinuities are relational or guarded.  A central-limit
statement becomes a finite estimate between rational laws plus explicit
truncation and smoothing errors.

\subsection{Topology}

\paragraph{Objects.}
Finite simplicial or cubical complexes, subdivision tables, chain complexes,
carrier clouds, finite covers, and quantitative compact metric presentations.

\paragraph{Primitive-recursive core.}
Boundary matrices, chain maps, parity, Sperner labels, simplicial subdivision,
finite homology certificates, and explicit prism homotopies.

\paragraph{Continuity interface.}
A map is observed through a modulus assigning a sufficiently fine source
level to each target level.  The finite map may be a cloud of target vertices.
A strong carrier condition makes every selection simplicial and any two
selections contiguous, so the induced homology map is independent of choice.
This resolves the noncanonical-selection issue without producing a global
continuous selector.

\subsection{Algebraic geometry}

Algebraic geometry can be split into exact finite algebra and genuinely
geometric producer theorems.

\paragraph{Affine presentations.}
An affine scheme of finite type is represented by a polynomial presentation
\(A=R[x_1,\ldots,x_n]/I\) together with bounded algebraic certificates.  A
morphism is a tuple of polynomials satisfying the target relations modulo
\(I\), witnessed by explicit combinations.  Equality of morphisms on a finite
presentation is another ideal-membership certificate.

\paragraph{Standard opens and gluing.}
A finite affine cover by standard opens \(D(f_i)\) carries an explicit cover
identity, for example \(\sum_i a_i f_i^N=1\) for a supplied common exponent
\(N\).  Local sections are fractions with synchronized
denominators.  Overlap compatibility is checked by cross multiplication after
a supplied localization exponent.  The finite affine-gluing theorem in
\cref{sec:casestudies} then constructs or verifies a global section.

\paragraph{Coherent sheaves.}
On a fixed finite affine cover, a finitely presented module is a matrix.  A
morphism is a matrix satisfying relation equations.  Kernels, cokernels, and
\v{C}ech complexes at a fixed degree are finite matrix calculations.  Exactness
is certified by explicit preimages or rank/minor data appropriate to the base
ring.

\paragraph{Points and dimensions.}
Bounded-height rational points are finite enumerations.  Dimension at a fixed
finite presentation can be certified through chains of prime ideals or finite
Gröbner data when an explicit degree bound is available.  General uniform
degree bounds and rational-point finiteness remain producer theorems, not
consequences of the interface alone.

\paragraph{Formal schemes and limits.}
An inverse system is represented only through finite truncations and explicit
transition maps.  A theorem about a requested truncation compiles directly.
A completed inverse-limit point requires a coherence/diameter theorem; it is
not obtained from levelwise nonemptiness alone.

\subsection{Category theory}

Finite categories, functors, natural transformations, adjunction data, and
finite diagrams are bounded tables.  Yoneda components, finite limits and
colimits, and finite Kan extensions reduce to table equations and bounded
search.  For categories of infinite mathematical structures, use categories
of finite presentations and contracted morphisms.  A categorical theorem is
then a reusable compiler for finite diagrams, not a quantification over an
uncoded universe of all objects.

\subsection{Differential and Riemannian geometry}

A charted object is represented by finitely many rational charts, overlap
margins, transition approximations, and derivative bounds.  Tensor identities
at rational sample points are finite algebra.  Curvature or geodesic
statements over a compact region require moduli and finite nets.  Coordinate
singularities are handled by relational chart clouds, while exact chart
selection requires a positive distance-to-boundary margin.

\subsection{Domain conversion table}

\begin{longtable}{@{}p{0.15\textwidth}p{0.23\textwidth}p{0.28\textwidth}p{0.22\textwidth}@{}}
\toprule
Domain & Finite presentation & Critical quantitative input & Safe output \\
\midrule
\endfirsthead
\toprule
Domain & Finite presentation & Critical quantitative input & Safe output \\
\midrule
\endhead
Algebra & coefficient/matrix tables & degree, pivot, or separation bound & exact certificate or cluster \\
Analysis & rational names and intervals & modulus, tail, apartness margin & regular name or guarded result \\
Diophantine geometry & bounded-height codes & height bound / covering representatives & complete finite list or conditional list \\
Analytic number theory & finite sums and kernels & effective cutoff and constants & explicit inequality \\
Probability & finite rational laws & tail and anti-concentration bounds & finite law approximation \\
Topology & finite complexes/carriers & subdivision and mesh modulus & chain/homology certificate \\
Algebraic geometry & finite presentations and affine covers & degree/localization/height bounds & finite algebra certificate or relational section \\
Category theory & finite categories/diagrams & finite bounds only & exact table theorem \\
Geometry & charts and rational nets & overlap and derivative margins & approximate invariant / guarded chart \\
\bottomrule
\end{longtable}

\section{Four developed case studies}
\label{sec:casestudies}

\subsection{Exact bisection for the inverse of \texorpdfstring{$y^3+y$}{y3+y}}

Let
\[
        f(y)=y^3+y.
\]
For rational input \(x\), we construct a regular rational name for the unique
solution of \(f(y)=x\).  The construction is useful because it exercises exact
rational arithmetic, a bounded loop, a monotonicity certificate, and a final
error contract without any transcendental machinery.

\begin{lemma}[Strict monotonicity]
For all rationals \(a<b\), \(f(a)<f(b)\).
\end{lemma}

\begin{proof}
Factor the difference:
\[
 f(b)-f(a)
 =(b-a)(a^2+ab+b^2+1).
\]
The first factor is positive.  The second is positive because
\[
 a^2+ab+b^2=\left(a+\frac b2\right)^2+\frac{3b^2}{4}\ge0.
\]
All identities are rational polynomial identities and are therefore suitable
for iPRA after denominator clearing.
\end{proof}

Set \(B=|x|+1\).  Then
\[
        f(-B)\le -B\le x\le B\le f(B),
\]
so \([-B,B]\) brackets the unique solution.  Define a bisection trace.  At
state \((\ell_i,h_i)\), let \(m_i=(\ell_i+h_i)/2\).  If
\(f(m_i)\le x\), set \((\ell_{i+1},h_{i+1})=(m_i,h_i)\); otherwise set it to
\((\ell_i,m_i)\).

\begin{lemma}[Bisection invariant]
For every \(i\),
\[
 f(\ell_i)\le x\le f(h_i),
 \qquad
 h_i-\ell_i=2B\cdot2^{-i}.
\]
\end{lemma}

\begin{proof}
The base case is the bracket above.  At a step, rational comparison decides
\(f(m_i)\le x\).  In the first branch, the lower inequality is updated to
\(f(m_i)\le x\), and the upper inequality is unchanged; the second branch is
symmetric.  In either branch the interval width is halved.  The conjunction is
a fixed Boolean combination of polynomial inequalities and rational equalities,
so it can be the trace invariant.
\end{proof}

Let \(q_p\) be the midpoint after \(N(p)\) steps, where \(N(p)\) is the
least integer satisfying
\[
       2B\cdot2^{-N(p)}\le 2^{-p-2}.
\]
Bounded doubling from the numerator and denominator of \(B\) computes this
integer primitively recursively.  Taking the least such integer makes \(N\)
nondecreasing, so the intervals used for successive precisions are nested.

\begin{theorem}[Cubic inverse name]
The sequence \((q_p)\) is a regular Cauchy name.  It determines the unique
real solution of \(y^3+y=x\), and every finite assertion needed for this claim
is provable from the bisection trace in iPRA.
\end{theorem}

\begin{proof}
The solution lies in every bracketing interval.  Hence the midpoint is within
half the interval width, at most \(2^{-p-3}\), of the solution.  For
\(r>p\), both \(q_p\) and \(q_r\) are within their respective half-widths of
the same nested solution interval, so
\[
       |q_r-q_p|\le 2^{-p-2}.
\]
This is stronger than the usual regularity requirement after a harmless index
shift.  Strict monotonicity gives uniqueness.  Internally, one need not
quantify over a completed real: the theorem states nested rational brackets,
shrinking widths, and the Cauchy inequality between midpoint terms.  These are
bounded rational calculations compiled by trace reflection.
\end{proof}

The included Python certificate for \(x=2\) and precision parameter \(12\)
records every rational interval and is independently checked by exact
cross-multiplication.

\subsection{Height bounds and complete lists}

Let \(P(c)\) be a decidable predicate on coded points, let \(h(c)\) be a
primitive-recursive height, and suppose there is a primitive-recursive finite
universe \(U(H)\) containing every well-formed code of height at most \(H\).

\begin{theorem}[Height bound implies complete list]
If
\[
        \forall c\,[P(c)\to h(c)\le H],
\]
then bounded filtering of \(U(H)\) by \(P\) produces a complete finite list of
all points satisfying \(P\).
\end{theorem}

\begin{proof}
Every retained code satisfies \(P\) by the filter definition.  Conversely, if
\(P(c)\), the height assumption puts \(c\) in the represented finite universe
\(U(H)\); the filter therefore retains it.  Duplicates may be removed by
bounded equality search.  All loops are bounded by the size of \(U(H)\).
\end{proof}

\begin{theorem}[Complete list implies height bound]
Suppose a finite list \(L\) is certified sound and complete:
\[
  c\in L\to P(c),
  \qquad
  P(c)\to c\in L.
\]
Then
\[
        H=\max\{h(c):c\in L\},
\]
with \(H=0\) for the empty list, is a height bound for every \(P\)-point.
\end{theorem}

\begin{proof}
If \(P(c)\), completeness gives \(c\in L\).  The defining finite-maximum
invariant says that the maximum is at least the height of every list entry.
Therefore \(h(c)\le H\).  The maximum is a primitive-recursive fold.
\end{proof}

The equivalence is often misunderstood.  It does \emph{not} say that a proof
of abstract finiteness yields a cutoff.  It says that either of two explicit
quantitative witnesses---a height bound or a certified complete list---can be
converted to the other by finite computation.

\subsection{Choice-independent simplicial maps from carriers}

Let \(K,L\) be finite simplicial complexes.  A strong carrier assigns to each
vertex \(v\) of \(K\) a nonempty finite set \(C(v)\) of vertices of \(L\), with
the property that for every simplex \(\sigma\in K\),
\[
      \bigcup_{v\in\sigma} C(v)
\]
is contained in a simplex of \(L\).  A selection chooses one
\(f(v)\in C(v)\) for each vertex.

\begin{theorem}[Carrier independence]
Every selection is a simplicial map \(f:K\to L\).  Any two selections
\(f,g\) are contiguous, and therefore induce the same homomorphism on homology
over \(\F_2\).
\end{theorem}

\begin{proof}
For a simplex \(\sigma\), the selected vertices \(f(v)\) lie in the union of
its carrier clouds, which lies in a target simplex.  Thus \(f\) is simplicial.
For two selections, the union
\(f(\sigma)\cup g(\sigma)\) is still contained in the same union of carrier
clouds, so it is a simplex; this is contiguity.

Define the prism operator on an ordered \(k\)-simplex by
\[
 P[v_0,\ldots,v_k]
 =\sum_{i=0}^k
 [f(v_0),\ldots,f(v_i),g(v_i),\ldots,g(v_k)]
 \quad\text{over }\F_2,
\]
omitting degenerate simplices according to the fixed normalization convention.
A finite cancellation calculation gives
\[
       \partial P+P\partial=f_\#+g_\#.
\]
Hence, for a cycle \(z\),
\(f_\#z+g_\#z=\partial Pz\), so the two images represent the same homology
class.  Every chain group and sum is finite; the prism identity is a bounded
matrix equality and can be certified directly.
\end{proof}

This theorem is the correct constructive replacement for choosing a canonical
simplicial approximation.  A finite carrier cloud is enough, because the
observable homology map is choice independent.

\subsection{Strong finite affine gluing}

Let \(R\) be a coded commutative ring and \(M\) a coded \(R\)-module.  Fix
\(f_0,\ldots,f_{r-1}\in R\) and an exponent \(N\), together with a cover
certificate
\begin{equation}
        \sum_{i<r} a_i f_i^N=1.
        \label{eq:coveridentity}
\end{equation}
A local section on \(D(f_i)\) is represented by \(m_i/f_i^N\).  Assume the
strong synchronized compatibility equations
\begin{equation}
        f_j^N m_i=f_i^N m_j
        \qquad(i,j<r).
        \label{eq:strongcompat}
\end{equation}
Define
\[
             m=\sum_{i<r}a_i m_i.
\]

\begin{theorem}[Finite affine gluing]
Under \eqref{eq:coveridentity} and \eqref{eq:strongcompat}, the global element
\(m\in M\) restricts to \(m_k/f_k^N\) on every \(D(f_k)\).  It is unique among
global elements whose differences are annihilated by a supplied common cover
power.
\end{theorem}

\begin{proof}
In the localization at \(f_k\), use the cover identity:
\[
 \frac{m_k}{f_k^N}
 =\sum_{i<r}a_i f_i^N\frac{m_k}{f_k^N}
 =\sum_{i<r}a_i m_i
 =m.
\]
The middle equality is exactly the localized form of
\eqref{eq:strongcompat}.  For uniqueness, if \(m'\) has the same restrictions
and a common exponent \(L\) satisfies
\(f_i^L(m-m')=0\) for every \(i\), use a certified identity
\(\sum b_i f_i^L=1\) to obtain
\[
 m-m'=\sum b_i f_i^L(m-m')=0.
\]
All sums, powers, and module equations are finite.
\end{proof}

Ordinary localization compatibility may require increasing denominators and
finding a common annihilating exponent.  The framework treats that exponent
and the associated cover identity as a finite producer certificate.  Once it
is supplied, gluing is a bounded algebraic trace.  This yields a concrete entry
point for affine schemes and coherent sheaves in iPRA without pretending that
arbitrary localization witnesses have a uniform bound for free.

\section{Resolution matrix for the reported difficulties}
\label{sec:difficulties}

The following matrix consolidates the recurring difficulties in the repository
and records the exact resolution supplied by the framework.  ``Resolved''
means that the generic mathematical argument is complete; it does not mean
that every domain-specific producer or every Lean proof has been completed.

\begin{longtable}{@{}p{0.14\textwidth}p{0.22\textwidth}p{0.31\textwidth}p{0.21\textwidth}@{}}
\toprule
Issue & Failure mode & Framework response & Status \\
\midrule
\endfirsthead
\toprule
Issue & Failure mode & Framework response & Status \\
\midrule
\endhead
Object-language gap & Host theorem or program is called an iPRA proof &
Require local \(\Prf\,\PRAAx\) proofs and compile them by finite-trace reflection &
Human proof complete; Lean milestones open \\

All-precision induction & Invariant contains \(\forall p\) &
Fix \(p\), build one finite query table, perform qf induction, quantify afterward &
Resolved schema \\

Bounded quantifiers & Induction predicate is not quantifier free &
Compile bounded \(\forall,\exists\) to Boolean folds and one equation &
Resolved; shared Lean fold library needed \\

Higher-type evaluator & Assumes one total evaluator for all PR codes &
Use fixed symbols, fixed families, or clocked traces &
Unrestricted evaluator impossible \\

Adaptive oracle use & Forward algorithm asks for unforeseen precision &
Run a reverse demand pass before the forward trace &
Resolved for static finite DAGs \\

Nonunique witnesses & Independent choices fail to cohere &
Finite candidate clouds, backward viability, stable quotient erasure &
Finite case resolved \\

Changing horizons & Least finite paths do not nest &
Diameter control plus motion-tail bound; no infinite branch selected &
Resolved conditional on modulus \\

Root and eigenvalue labels & Labels jump at collisions &
Return clusters/multisets; exact labels only under separation &
Representation repaired \\

Puiseux branches & Exact coefficients requested from arbitrary Cauchy input &
Use exact algebraic splitting data or approximate relational branches &
Unrestricted producer remains obstructed \\

Rotation axis & Exact unit axis requested at identity &
Use quaternion/projective cloud; exact axis only with gap from identity &
Representation repaired \\

Partial operations & Division or chart selection from mere \(x\ne0\) &
Totalize and consume a positive rational apartness margin &
Resolved under guard \\

Compactness & Abstract subcover or minimizer &
Finite net/cover producer, coverage certificate, modulus, bounded minimization &
Resolved for effective presentations \\

Metastability & Stable windows treated as a Cauchy rate &
Retain window theorem; require an explicit tail modulus for a name &
General conversion impossible \\

Diophantine finiteness & Cardinality mistaken for a search bound &
Separate height/list producer from generic enumeration and verification &
Verifier resolved; producer theorem-specific \\

Mordell--Weil descent & Selmer class mistaken for a rational covering point &
Require certified rational representatives or a relational descent object &
Producer open in general \\

Ineffective constants & ``Sufficiently large'' hides noncomputable data &
Make every cutoff and quantifier dependence explicit &
Conditional verifier resolved \\

Probability & General measures and discontinuous CDFs hide data &
Finite rational laws plus tail, quantization, and anti-concentration moduli &
Resolved for effective laws \\

Topology & Compact spaces used without numerical control &
Finite complexes, carriers, subdivisions, parity and chain certificates &
Representation resolved; theorem invariant remains \\

Algebraic geometry & Localizations and sheaves treated abstractly &
Finite affine presentations, synchronized denominators, cover and gluing certificates &
Strong finite gluing resolved; bounds conditional \\

Exten\-sion\-al\-ity & Code equality confused with object equality &
Observational setoids and contracted congruence proofs &
Resolved \\

Lexicographic descent & Termination asserted without a loop bound &
Compile a rank bound and stuttering terminal transition to a bounded trace &
Resolved with rank producer \\

Classical case split & Exact sign/side selected from approximate data &
Three-way guard or relational union; collapse under apartness &
Unrestricted split implies LLPO-type principle \\
\bottomrule
\end{longtable}

\subsection{Priority ordering for future work}

The highest-value engineering milestone is the generic trace compiler, because
it turns several existing source-level constructions into actual iPRA proofs.
The next milestone is a shared bounded-Boolean library.  Third is the object-
language viability dynamic program.  Only after these common layers should
large domain packets be migrated.  Deep producer theorems should be tackled in
parallel but kept in separate modules so that they do not block verification of
the generic compiler.

\subsection{Acceptance criterion for a migrated packet}

A packet should be considered fully embedded only when it supplies:
\begin{enumerate}
\item a pinned source audit and exact target statement;
\item fixed \(\PR\) symbols or clocked evaluators for every algorithm;
\item explicit moduli, guards, bounds, and finite universes;
\item a theorem ledger distinguishing producers, verifiers, and barriers;
\item local \(\Prf\,\PRAAx\) derivations;
\item a generic compiler invocation yielding the main derivation;
\item a build with no \code{sorry}, \code{axiom}, or untracked external theorem;
\item optional host-level tests clearly marked as regression evidence only.
\end{enumerate}

\section{Lean formalization plan}
\label{sec:leanplan}

\subsection{Pinned upstream API}

The outline imports only checked repository modules and uses their public
objects:
\[
\code{PR},\quad \code{PTm},\quad \code{PRAAx},\quad
\code{Prf},\quad \code{PRA.qfInd},\quad
\code{NFm.ofPRA},\quad \code{pra_iff_native}.
\]
The essential proof combinators are
\code{Prf.subst}, \code{Prf.weaken}, \code{Prf.trans},
\code{Prf.impI}, \code{Prf.impE}, \code{Prf.allI},
\code{Prf.allE}, \code{Prf.eqSubst1}, \code{Prf.eqSymm}, and
\code{Prf.eqCongr}.  Primitive-recursion equations come from
\code{PRA.precZero} and \code{PRA.precSucc}.

\subsection{File dependency graph}

\begin{enumerate}
\item \code{ObjectTraceCompiler.lean}: formula builders, trace system, run
equations, four proof milestones, main compiler.
\item \code{BoundedFormulaCompiler.lean}: Boolean symbols, folds, fixed bounded
formula syntax, reflection theorem.
\item \code{RelationalWitnesses.lean}: finite tower codes, viability, least
path, stable observation.
\item \code{CertifiedProgram.lean}: finite program/DAG syntax, reverse demands,
forward compilation to a trace system.
\item \code{ApplicationTargets.lean}: guarded reciprocal, height/list, cubic
bisection, carrier, and affine-gluing targets.
\end{enumerate}

The first file can be completed independently and should be merged before the
others.  It is the narrowest bridge to the checked backend.

\subsection{Exact trace layouts}

The object trace file uses these de Bruijn layouts:
\begin{center}
\begin{tabular}{lll}
\toprule
Formula & Arity & Variables by index \\
\midrule
\code{BaseFm} & 1 & \([u]\) \\
\code{StepFm} & 3 & \([i,s,u]\) \\
\code{FinalFm} & 2 & \([s,u]\) \\
\code{InvRunFm} & 2 & \([i,u]\) \\
\code{ResultFm} & 0 & closed, with one universal binder \\
\bottomrule
\end{tabular}
\end{center}
Term builders \code{app1}, \code{app2}, and \code{app3} make argument order
visible and reduce substitution errors.

\subsection{Milestone 1: base case}

The target is
\begin{lstlisting}[language=Lean]
theorem trace_base (T : PRTraceSystem) (h : TraceProofs T) :
  [GoodFm T] |- [PRAAx] (InvRunFm T).subst (inst zero)
\end{lstlisting}
The proof steps are:
\begin{enumerate}
\item instantiate \code{h.base} at \code{var 0};
\item weaken it to context \code{[GoodFm T]};
\item eliminate the implication using \code{Prf.hyp0};
\item obtain \code{run_zero T (var 0)};
\item transport the state argument with \code{Prf.eqSubst1};
\item close substitution equalities by \code{simp} lemmas for \code{app1/2/3}.
\end{enumerate}
A useful helper is a congruence lemma for \code{IsTrue (app3 T.inv i - u)}.

\subsection{Milestone 2: successor case}

The target is exactly the successor premise expected by \code{PRA.qfInd}:
\begin{lstlisting}[language=Lean]
theorem trace_step (T : PRTraceSystem) (h : TraceProofs T) :
  InvRunFm T :: [GoodFm T].map Fm.shift |- [PRAAx]
    (InvRunFm T).subst (repl (succ (var 0)))
\end{lstlisting}
Instantiate \code{h.step} with
\((i,\Run(i,u),u)\), eliminate implications using \code{Prf.hyp0} and
\code{Prf.hyp1}, and transport along \code{run_succ}.  Most work is normalizing
the shifted guard and the simultaneous substitution.

\subsection{Milestone 3: induction and final contract}

The induction theorem is already one line:
\begin{lstlisting}[language=Lean]
exact PRA.qfInd (InvRunFm T) (invRun_qf T)
  (trace_base T h) (trace_step T h)
\end{lstlisting}
Then use
\begin{lstlisting}[language=Lean]
Prf.allE (app1 T.bound (var 0)) (trace_all T h)
\end{lstlisting}
to obtain the invariant at the terminal index.  Instantiate \code{h.final} with
the terminal run and input, eliminate implications, discharge the guard, and
close the input binder.

\subsection{Milestone 4: bounded Boolean folds}

Introduce canonical Boolean-valued symbols and equations.  For each fixed body
symbol \(f\), define \code{allFold f}, \code{anyFold f}, and
\code{leastFold f} by \code{PR.prec}.  Prove:
\begin{itemize}
\item zero and successor equations from \code{PRA.precZero/precSucc};
\item Boolean range invariants;
\item reflection against bounded formula abbreviations;
\item least-witness soundness and minimality.
\end{itemize}
Use a fixed family indexed at the Lean level; do not define a universal PR
evaluator internally.

\subsection{Milestone 5: relational viability}

Encode a tower by fixed symbols for level bounds, candidate validity, and
edges.  Define viability by a backward recursion on the distance to horizon.
The principal object theorem is a Boolean equivalence between viability and a
bounded path verifier.  A second theorem proves correctness of least viable
successor.  Both use quantifier-free induction on the finite distance.

\subsection{Milestone 6: certified program compilation}

Replace semantic \code{Prop} fields in the program outline by:
\begin{itemize}
\item formula builders over \code{PTm};
\item local derivation fields;
\item fixed arity and state-layout records;
\item a compiler producing \code{PRTraceSystem} and \code{TraceProofs}.
\end{itemize}
Sequential composition uses \code{Prf.trans}.  Bounded loops invoke the trace
compiler recursively.  A static DAG is first lowered to reverse-demand and
forward-table loops, then compiled once.

\subsection{Suggested build discipline}

\begin{enumerate}
\item Place the modules under the root Lake workspace, in namespace
\code{LeanProofs.ConciseIPRA.}\allowbreak\code{FinitaryInterface}.
\item Add a lint command rejecting \code{sorryAx}, explicit \code{axiom}, and
unintended imports.
\item Build after each milestone; keep formula-normalization lemmas separate
from mathematical proofs.
\item Add tiny closed examples: identity trace, addition fold, least witness,
three-level viability tower, and cubic bisection step.
\item Only after the generic modules are checked should large packets import
them.
\end{enumerate}

\subsection{Estimated residual risk}

The trace compiler is low mathematical risk and moderate de Bruijn-engineering
risk.  Boolean folds are moderate risk because a clean fixed-family API must be
chosen.  Relational viability is low mathematical risk but may produce large
formula terms.  The program compiler is architectural rather than logical;
its main risk is over-generalization.  It should initially support only static
DAGs and bounded loops, adding richer syntax after concrete migrations expose
a need.

\section{Barriers, exact status, and conclusions}
\label{sec:barriers}

\subsection{No total primitive-recursive universal evaluator}

\begin{proposition}
There is no primitive-recursive function \(U(e,x)\) that enumerates all unary
primitive-recursive functions, in the sense that for every unary PR function
\(f\) there is an \(e\) with \(U(e,x)=f(x)\) for all \(x\).
\end{proposition}

\begin{proof}
Assume such a \(U\).  Then
\[
       d(x)=U(x,x)+1
\]
is primitive recursive.  By universality there is an index \(e\) with
\(U(e,x)=d(x)\) for all \(x\).  At \(x=e\),
\[
       U(e,e)=d(e)=U(e,e)+1,
\]
a contradiction.
\end{proof}

Thus a generic compiler may be parametric over fixed \(\PR\) syntax at the
Lean level, but it may not assume an internal total evaluator for arbitrary PR
codes.

\subsection{Finiteness does not uniformly yield a cutoff}

\begin{proposition}
There is no computable operation that, from every decidable predicate known to
have at most one witness, returns a bound containing that witness whenever it
exists.
\end{proposition}

\begin{proof}
For a machine code \(e\), let \(P_e(n)\) say that machine \(e\) halts for the
first time at exactly step \(n\).  This is a bounded-step decidable predicate
and has at most one witness.  If a uniform bound \(H(e)\) existed, simulate the
machine for \(H(e)\) steps.  It halts eventually if and only if it halts within
that bound, deciding the halting problem.
\end{proof}

This is why abstract cardinality bounds in Diophantine geometry cannot be fed
directly to bounded search.  A height bound or complete list is genuinely
additional data.

\subsection{Exact sign selection entails a constructive omniscience principle}

Consider a binary sequence \(a_n\) with at most one entry equal to one, and the
regular real
\[
 x=\sum_{n\ge0}(a_{2n}-a_{2n+1})2^{-n-1}.
\]
If the unique one occurs at an even index, \(x>0\); if it occurs at an odd
index, \(x<0\); if no one occurs, \(x=0\).  A uniform operation that, for every
real name, returned a justified side \(x\ge0\) or \(x\le0\) would decide the
LLPO disjunction
\[
   (\forall n\;a_{2n}=0)\quad\vee\quad
   (\forall n\;a_{2n+1}=0).
\]
Such a selector is not available in the intended intuitionistic setting.
Three-way guards and relational branch unions avoid this principle; an
apartness margin makes exact selection constructive.

\subsection{Metastability does not uniformly give a Cauchy modulus}

A bounded increasing computable rational sequence can encode halting
information.  Let \(q_s\) add weight \(2^{-e-1}\) when machine \(e\) is seen to
halt by stage \(s\).  The sequence is increasing and bounded by one.  At each
fixed precision it has effective finite metastability bounds based on the fact
that there can be only finitely many jumps larger than that precision.  A
uniform conversion of such metastability data into a tail modulus would decide,
for each finite set of machines, whether any new halt can occur after the
claimed modulus stage, and hence decide halting.  Therefore a finite-window
theorem and a Cauchy-name theorem must remain different interfaces.

\subsection{Branch data from arbitrary Cauchy coefficients}

Exact Puiseux or eigenvector branches can encode discontinuous choices at
singular or multiple points.  A uniform exact branch producer from arbitrary
coefficient names would in particular decide whether a small perturbation lies
on one side of a discriminant locus.  The safe interfaces are exact algebraic
input with splitting data, approximate branch clouds, or exact branches guarded
by a separation certificate.

\subsection{Status of this package}

\begin{longtable}{p{0.26\textwidth}p{0.24\textwidth}p{0.42\textwidth}}
\toprule
Component & Status & Evidence \\
\midrule
\endfirsthead
\toprule
Component & Status & Evidence \\
\midrule
\endhead
Human framework theorems & complete at the stated level & this report, with explicit hypotheses and proofs \\
Python reference kernel & executed & eleven exact-arithmetic regression tests \\
Serialized example & verified & cubic inverse certificate for \(x=2\) \\
Lean formula definitions & source outline & exact pinned API and documented arities \\
Generic trace proof & mathematical proof complete; Lean incomplete & four small visible \code{sorry} milestones \\
Boolean compiler & mathematical proof complete; Lean incomplete & fixed-family design and fold plan \\
Relational viability & mathematical proof complete; Lean incomplete & exact Python dynamic program and Lean targets \\
Bi-interpretation transport & backend-checked upstream & \code{pra_iff_native} in \code{ConciseIPRA} \\
General height/constant producers & not claimed & isolated as theorem-specific obligations \\
\bottomrule
\end{longtable}

\subsection{Main conclusion}

The broad range of mathematics represented in the repository does not require
a separate foundational trick for every theorem.  A stable common core has
emerged:
\[
\boxed{\begin{gathered}
\text{finite observations}+\text{backward demands}+\text{bounded certificates}\\
{}+\text{relational coherence}+\text{object-language trace reflection}
\end{gathered}}
\]
This core is strong enough to organize algebra, analysis, number theory,
topology, probability, category theory, and finite parts of algebraic geometry.
It resolves the recurrent logical difficulties without overstating deep
mathematical results.  Most importantly, it identifies a concrete next step:
finish and kernel-check the small generic \code{traceReflection} theorem, then
migrate domain packets through a shared Boolean and viability library.  That
work would convert a substantial body of existing finitary mathematics from
plausible formalization plans into actual iPRA derivations.


\appendix
\section{Theorem summary and dependency map}

\begin{longtable}{p{0.11\textwidth}p{0.48\textwidth}p{0.31\textwidth}}
\toprule
ID & Theorem & Depends on \\
\midrule
\endfirsthead
\toprule
ID & Theorem & Depends on \\
\midrule
\endhead
T1 & Bounded fold reflection & PR recursion equations, qf induction \\
T2 & Fixed bounded formula Booleanization & T1, Boolean truth tables \\
T3 & Fixed-precision induction & T2, universal introduction \\
T4 & Finite-trace reflection & run equations, qf induction, equality transport \\
T5 & Sequential contract composition & substitution and simultaneous cut \\
T6 & Reverse-demand soundness & finite topological induction, restriction maps \\
T7 & Certified DAG compilation & T2, T4, T6 \\
T8 & Viability/path equivalence & bounded backward induction \\
T9 & Least viable finite path & T1, T8 \\
T10 & Stable quotient erasure & T8--T9, finite stability premise \\
T11 & Shrinking viability gives a regular name & T8, diameter and tail bounds \\
T12 & Guarded reciprocal & rational interval inequalities, apartness margin \\
T13 & Finite-net minimization & bounded minimum, modulus \\
T14 & Height bound implies complete list & bounded universe and filter \\
T15 & Complete list implies height bound & finite maximum \\
T16 & Cubic inverse regular name & trace reflection, monotonicity, bisection \\
T17 & Carrier choice independence & finite contiguity prism identity \\
T18 & Strong finite affine gluing & cover identity, synchronized compatibility \\
\bottomrule
\end{longtable}

\subsection{Backend dependency}

T1--T18 are intended to live in ordinary iPRA.  Their final checked transport
to the native scalar calculus is not reproved here; it is supplied by the
upstream theorem \code{pra_iff_native}.  Domain packets should depend on the
ordinary compiler layer and invoke the native bridge only at their public
boundary.

\section{Packet conversion checklist}

For a theorem packet, complete the following worksheet.

\begin{enumerate}
\item \textbf{Target statement.}  Rewrite every asymptotic, limiting, or
existence assertion with explicit numerical quantifier order.
\item \textbf{Input presentation.}  List all finite codes, well-formedness
predicates, and observation queries.
\item \textbf{Producers.}  Identify every modulus, height bound, cover,
separation margin, representative, or analytic constant.  Mark whether it is
proved, assumed conditionally, or obstructed.
\item \textbf{Finite verifier.}  Give a bounded Boolean predicate for each
certificate.
\item \textbf{Partial operations.}  Totalize them and add positive rational
guards.
\item \textbf{Nonuniqueness.}  Decide whether to use exact canonicalization, a
finite relation, or stable quotient erasure.
\item \textbf{Precision flow.}  Define a primitive-recursive backward demand
for every construction and join fan-out demands.
\item \textbf{Bounded formula compilation.}  Replace every bounded quantifier
inside a loop invariant by a fold.
\item \textbf{Trace.}  Specify state, initialization, bound, step, invariant,
and acceptance.
\item \textbf{Local proofs.}  Produce closed \(\Prf\,\PRAAx\) derivations of
base, preservation, and final soundness.
\item \textbf{Compiler invocation.}  Apply trace reflection and contract
composition.
\item \textbf{Extensionality.}  Prove that observationally equivalent inputs
give equivalent outputs.
\item \textbf{Lean audit.}  Build with no \code{sorry} or unintended axioms;
record the exact repository commit.
\item \textbf{Evidence separation.}  Keep Python tests and external
certificates as regression evidence, not theorem status.
\end{enumerate}

A packet failing at item 3 has a mathematical producer problem.  A packet
failing at item 6 may have an ill-posed output.  A packet reaching item 10 but
not item 11 has the generic compiler gap addressed by this report.

\section{Pinned \texorpdfstring{\code{ConciseIPRA}}{ConciseIPRA} API inventory}

This appendix records the upstream names used by the Lean outline.  It is not a
substitute for the source files; it is a migration checklist.

\begin{longtable}{p{0.29\textwidth}p{0.62\textwidth}}
\toprule
Name & Role \\
\midrule
\endfirsthead
\toprule
Name & Role \\
\midrule
\endhead
\code{PR k} & primitive-recursive function symbols of arity \(k\) \\
\code{PTm n} & terms with \(n\) free variables \\
\code{PTm.app}, \code{PTm.zero}, \code{PTm.succ} & term constructors \\
\code{Fm PTm n} & first-order formulas \\
\code{Fm.QF} & quantifier-free syntactic predicate \\
\code{PRAAx} & iPRA axiom family \\
\code{Prf PRAAx Gamma phi} & object-language derivability \\
\code{PRA.precZero} & primitive recursion zero equation at arbitrary terms \\
\code{PRA.precSucc} & primitive recursion successor equation at arbitrary terms \\
\code{PRA.qfInd} & quantifier-free induction rule \\
\code{Prf.subst} & simultaneous substitution in a derivation \\
\code{Prf.weaken} & context weakening \\
\code{Prf.trans} & simultaneous cut / replacement of hypotheses \\
\code{Prf.impI}, \code{Prf.impE} & implication rules \\
\code{Prf.allI}, \code{Prf.allE} & universal quantifier rules \\
\code{Prf.eqSubst1} & one-variable Leibniz transport \\
\code{Prf.eqSymm}, \code{Prf.eqTrans} & equality symmetry and transitivity \\
\code{Prf.eqCongr}, \code{Prf.eqCongr1} & term congruence \\
\code{NFm.ofPRA} & translation from ordinary iPRA formula to native formula \\
\code{pra_iff_native} & equivalence of ordinary and native derivability \\
\bottomrule
\end{longtable}

The outline intentionally avoids private implementation details of the native
calculus.  It constructs ordinary iPRA proofs first, so future changes to the
compact presentation need affect only the already checked bridge.


\begin{thebibliography}{20}
\bibitem{ProveItRoot}
V. Reshetnikov, \emph{ProveIt}, repository snapshot \RepoCommit,
\url{https://github.com/VladimirReshetnikov/ProveIt}.

\bibitem{ConciseIPRA}
V. Reshetnikov, \emph{Concise presentations of intuitionistic primitive recursive arithmetic},
\code{Logic/PrimitiveRecursiveArithmetic/ConciseIPRA} in \cite{ProveItRoot}.

\bibitem{PacketsIndex}
V. Reshetnikov, \emph{Delivered iPRA proof packets},
\path{Logic/PrimitiveRecursiveArithmetic/docs/proofs/README.md} in \cite{ProveItRoot}.

\bibitem{BishopBridges}
E. Bishop and D. Bridges, \emph{Constructive Analysis}, Springer, 1985.

\bibitem{BridgesRichman}
D. Bridges and F. Richman, \emph{Varieties of Constructive Mathematics}, Cambridge University Press, 1987.

\bibitem{TroelstraVanDalen}
A. S. Troelstra and D. van Dalen, \emph{Constructivism in Mathematics}, Volumes I--II, North-Holland, 1988.

\bibitem{Kohlenbach}
U. Kohlenbach, \emph{Applied Proof Theory: Proof Interpretations and their Use in Mathematics}, Springer, 2008.

\bibitem{Weihrauch}
K. Weihrauch, \emph{Computable Analysis}, Springer, 2000.
\end{thebibliography}

\end{document}
